\documentclass{book}

\usepackage[utf8]{inputenc}  % para que funcionen las tildes
\usepackage{amsmath}
\usepackage{amssymb}
\usepackage{amsthm}
\usepackage{letltxmacro}
\LetLtxMacro\amsproof\proof
\LetLtxMacro\amsendproof\endproof
\usepackage{txfonts}
\usepackage{lmodern}
\usepackage[dvipsnames]{xcolor}
\usepackage{thmtools, thm-restate}
\usepackage{datetime} % para la hora de compilación
\usepackage{graphicx}
\graphicspath{{images/}}
\usepackage[spanish,es-noquoting]{babel} % es-noquoting es para que funcione tikz
\usepackage{mathabx} % para \divides
\usepackage{centernot} % para \centernot\inculdes
\usepackage{wrapfig}
\usepackage{multicol}
\usepackage{subcaption}
\usepackage{etoolbox}
\usepackage{xfrac}
\usepackage{glossaries}
\usepackage{standalone}
\usepackage{pgfmath}
\usepackage{pgfplots}
\usepackage{pdfpages}
\usepackage{tikz}
\usepackage{mathtools}
\usepackage{tikz-qtree}
\usepackage{cancel}
\usepackage[shortlabels]{enumitem}

\usetikzlibrary{arrows.meta}
\usetikzlibrary{shapes}
\usetikzlibrary{decorations.text}
\usetikzlibrary{positioning}
\usetikzlibrary{external}
\tikzexternalize[prefix=./tikzbuild/]
\tikzexternalize % activate!

\AtBeginDocument{%
  \LetLtxMacro\proof\amsproof
  \LetLtxMacro\endproof\amsendproof
}

\definecolor{coolblack}{rgb}{0.0, 0.18, 0.39}
%\usepackage[a6paper,margin=5mm]{geometry}
\usepackage[a4paper, top=2.5cm,bottom=2.5cm,left=2.5cm,right=2.5cm]{geometry}

%\setlength{\parindent}{0pt}
\usepackage{parskip}

\usepackage{hyperref}

% ESTILOS DE DEFINICIONES Y TEOREMAS
\declaretheoremstyle[
	bodyfont=\normalfont,
	shaded={
		margin=8pt,
		bgcolor=White,
		rulecolor=Black,
		rulewidth=1pt
}]{mythm}

\declaretheoremstyle[
	bodyfont=\normalfont,
	shaded={
		margin=1em,
		bgcolor={rgb}{0.9,0.9,0.9}
}]{mydfn}

\declaretheoremstyle[
bodyfont=\normalfont,
spacebelow=1em,
spaceabove=1em,
]{myej}

\declaretheoremstyle[
bodyfont=\normalfont,
spacebelow=1em,
spaceabove=1em,
shaded={
    margin=8pt,
    bgcolor=White,
    rulecolor=coolblack,
    rulewidth=1pt
}]{thej}

\declaretheoremstyle[
bodyfont=\normalfont,
spacebelow=1em,
spaceabove=1em,
thmbox=M
]{myeg}

% DEFINICIONES DE ENTORNOS DE TEOREMAS
\declaretheorem[
	name=Teorema,
	refname={teorema,teoremas},
	Refname={Teorema,Teoremas},
	style=mythm
]{thm}

\declaretheorem[
	name=Corolario,
	refname={corolario,corolarios},
	Refname={Corolario,Corolarios},
	style=myej
]{cor}

\declaretheorem[
	name=Proposici\'{o}n,
	refname={proposici\'{o}n,proposiciones},
	Refname={Proposici\'{o}n, Proposiciones},
	sharenumber=thm,
	style=myej
]{pro}

\declaretheorem[
name=Lema,
refname={lema,lemas},
Refname={Lema, Lemas},
sharenumber=thm,
style=myej
]{lm}

%\newtheorem*{cor}{Corolario}
\newtheorem*{lem}{Lema}

\declaretheorem[
	name=Definici\'{o}n,
	refname={definici\'{o}n,definiciones},
	Refname={Definici\'{o}n,Definiciones},
	style=mydfn
]{dfn}

\declaretheorem[
name=Ejercicio,
refname={ejercicio,ejercicios},
Refname={Ejercicio,Ejercicios},
style=myej,
numbered=no
]{ex}

\declaretheorem[
name=Ejercicio propuesto,
refname={ejercicio,ejercicios},
Refname={Ejercicio,Ejercicios},
style=thej,
]{th_ex}

\declaretheorem[
name=Ejemplo,
refname={ejemplo,ejemplos},
Refname={Ejemplo,Ejemplos},
style=myeg
]{eg}

\declaretheorem[
name=Observaci\'{o}n,
refname={observaci\'{o}n,observaciones},
Refname={Observaci\'{o}n,Observaciones},
style=myej,
numbered=no
]{obs}

\declaretheorem[
name=Notaci\'{o}n,
refname={notaci\'{o}n},
Refname={Notaci\'{o}n},
style=myej,
numbered=no
]{nota}

\makeglossaries
\newglossaryentry{di}{name={DI},description={Dominio de integridad (ver \autoref{dfn:dominiointegridad})}}
\newglossaryentry{dip}{name={DIP},description={Dominio de ideales principales (ver \autoref{dfn:dominioidealesprincipales})}}
\newglossaryentry{dfu}{name={DFU},description={Dominio de factorización única (ver \autoref{dfn:dfu})}}
\newglossaryentry{de}{name={DE},description={Dominio euclídeo (ver \autoref{dfn:de})}}


\setcounter{chapter}{0}
%COMANDOS ÚTILES PARA DERIVADAS
\AtBeginDocument{\renewcommand{\d}{\mathrm{d}}}
\newcommand{\Dd}[1]{\frac{\mathrm{d}}{\mathrm{d}#1}}
\newcommand{\dd}[2]{\frac{\mathrm{d}#1}{\mathrm{d}#2}}
\newcommand{\mbf}[1]{\mathbf{#1}}
\newcommand{\mtc}[1]{\mathcal{#1}}
% COMANDOS ÚTILES PARA LA TEORÍA DE GRUPOS
%\newcommand{\normsub}{\mathbin{\triangleleft}}
\newcommand{\normsub}{\lhd}
\newcommand{\uds}[1]{\mathcal{U}(#1)}
\newcommand{\inv}[1]{#1^{-1}}
\newcommand{\ima}{\text{Im }}
\newcommand{\isom}{\simeq}
\newcommand{\autom}[1]{\text{Aut}(#1)}
\newcommand{\gen}[1]{\langle#1\rangle}
\newcommand{\biy}[1]{\text{Biy}(#1)}

%\DeclareMathSymbol{\varprod}{\mathop}{largesymbolsA}{16}

\newcommand{\hr}{\rule{\textwidth}{.4pt}}
\newcommand{\N}{\mathbb{N}}
\newcommand{\Z}{\mathbb{Z}}
\newcommand{\R}{\mathbb{R}}
\newcommand{\Q}{\mathbb{Q}}
\newcommand{\C}{\mathbb{C}}
\newcommand{\A}{\mathbb{A}}
\newcommand{\F}{\mathbb{F}}
\newcommand{\J}{\mathbb{J}}
\newcommand{\B}{\mathbb{B}}
\DeclareMathOperator{\diag}{diag}
\DeclareMathOperator{\traza}{traza}
\renewcommand{\P}{\mathcal{P}}
\renewcommand{\S}{\mathbb{S}}
\newcommand{\ZnZ}{\mathbb{Z}/n\mathbb{Z}}
\newcommand{\ZmZ}{\mathbb{Z}/m\mathbb{Z}}
\newcommand{\0}{\mathbf{0}}
\newcommand{\1}{\mathbf{1}}
\newcommand{\gr}{\text{gr }}
\newcommand\ninf[1]{\left|\left|#1\right|\right|_\infty}
\newcommand{\sdf}[1]{\left\{#1\right\}}

%% LOGICA
\newcommand{\bicond}{\leftrightarrow}
\newcommand{\then}{\rightarrow}
\newcommand{\ter}{\mathrm{Ter}}
\newcommand{\for}{\mathrm{For}}
\newcommand{\sbn}{\mathrm{SBN}}
\newcommand{\model}{\mathrm{Mod}}
\newcommand{\te}{\mathrm{te}}
\newcommand{\quo}[2]{{\left.\raisebox{.2em}{$#1$}\middle/\raisebox{-.2em}{$#2$}\right.}}
\newcommand{\quox}[1]{\quo{#1}{x}}
\newcommand{\implied}{\ \Longleftarrow\ }
%%%% MATRICES
\newcommand\mx[1]{\begin{bmatrix*}#1\end{bmatrix*}}

\newcommand{\suc}[3][1]{{#2}_{#1}, \ldots, {#2}_{#3}}
\newcommand{\sucn}[1]{\suc{#1}{n}}

\renewcommand\qedsymbol{$\diamondsuit$}
\setcounter{MaxMatrixCols}{20}

\AtBeginEnvironment{align}{\setcounter{equation}{0}}
\renewcommand{\theequation}{\arabic{equation}}

\title{Apuntes de L\'ogica M\'atematica}
\author{Rafael S\'{a}nchez S\'{a}nchez}
\begin{document}
\begin{titlepage}
	\includepdf{portada.pdf}
\end{titlepage}


Revisión del \today $ $ a las \currenttime.

\begin{center}

\end{center}

\tableofcontents
\part{Sintaxis y sem\'antica de primer orden}
% !TeX root = ../logica-matematica.tex

\chapter{Estructuras y lenguajes de primer orden}
\section{Introducción}
\begin{dfn}[Estructura]\label{dfn:estructura}
    Una \textbf{estructura} (de primer orden) $\mtc{A}$ consta de un conjunto no vacío $A$ (universo) y un conjunto de funciones, relaciones y elementos del universo.
\end{dfn}
Puede parecer una definición algo abstracta, así que vamos a ver algunos ejemplos:
\begin{eg}[Estructuras. Ejemplos]
    \begin{itemize}
        \item $ \mtc{G} = \gen{G, \cdot, 1} $
        \item $ R = \gen{\R,\  +,\  -,\  \cdot,\  0,\  1,\  \leq} $
        \item $ R_* = \gen{\R,\  +,\  -,\  \cdot,\  0,\  1,\ \left. \sin \right|_{[0, 1]},\  \exp}$
    \end{itemize}
\end{eg}

\begin{dfn}[Lenguaje]\label{dfn:lenguaje}
    Un \textbf{lenguaje} (de primer orden) consta de:\\
    \begin{itemize}
        \item Para cada $n \in \N^*$ un conjunto $\mtc{F}_n$ de símbolos de funciones n-arias.
        \item Para cada $m \in \N^*$ un conjunto $\mtc{R}_m$ de símbolos de relaciones m-arias.
        \item Un conjunto $\mtc{C}$ de constantes.
    \end{itemize}
    Además de conjuntos de símbolos lógicos, variables $\{v_i\}_{i\in \N}$, cuantificadores $(\exists,\ \forall)$, conectores $(\land, \lor, \iff, \implies)$, símbolos de igualdad $=$, y paréntesis $()$. No se suelen especificar en la declaración del lenguaje.
\end{dfn}

\begin{eg}[Lenguajes. Ejemplos]\label{eg:lenguajes}
    \begin{itemize}
        \item $L = \left\{\ast,\ e\right\}$
        \item $L = \left\{\cdot,\ 1\right\}$
        \item $L = \left\{+,\ 0\right\}$
        \item $L = \left\{\cdot,\ ^{-1},\ 1\right\}$
    \end{itemize}
    Además, para hablar del comportamiento de los reales vamos a usar el lenguaje:
    $$
        L = \left\{+,\ -,\ \cdot,\ 0,\ 1,\ \leq \right\}
    $$
    Es importante destacar que cuando escribimos $+$ en la declaración del lenguaje no nos referimos a la función: $+: \R^2 \to \R$ si no a un símbolo que luego interpretamos como tal función en la estructura.
\end{eg}

\section{L-estructuras}
\begin{dfn}[L-estructura]
    Dado un lenguaje \textsuperscript{\ref{dfn:lenguaje}} $L$, una \textbf{L-estructura} $\mtc{A}$ o una \textbf{interpretaci\'on de L} consta de:
    \begin{itemize}
        \item un universo $A \neq \varnothing$
        \item una función n-aria $f^\mtc{A} : A^n \to A$ para cada símbolo de función $f \in \mtc{F}_n$
        \item una relación m-aria $R^\mtc{A} \subseteq A^m$ para cada símbolo de relación $R \in \mtc{R}_m$
        \item un elemento $e^\mtc{A}\in A$ para cada constante $c \in \mtc{C}$.
    \end{itemize}
\end{dfn}

\begin{nota}
    En la definición anterior:\\
    $L = \mtc{F} \cup \mtc{R} \cup \mtc{C}$ (lenguaje), $\mtc{F}=\bigcup_{n\in \N^*} \mtc{F}_n$ (símbolos de función), $\mtc{R}=\bigcup_{m\in \N^*} \mtc{R}_m$ (símbolos de relación), $\mtc{C}$ constantes.\\
    Además, para cada símbolo $s\in L$, $s^\mtc{A}$ es la interpretación de $s$ a $\mtc{A}$.
\end{nota}

\begin{eg}[L-estructura. Ejemplos]
    \begin{itemize}
        \item $R = \gen{\R,\ +^R,\ -^R,\ \cdot^R,\ \leq^R,\ 0^R,\ 1^R}$ (interpretación del lenguaje de los reales con el universo $\R$).
        \item $\mtc{A} = \gen{A,\ +^\mtc{A},\ -^\mtc{A},\ \cdot^\mtc{A},\ \leq^\mtc{A},\ 0^\mtc{A},\ 1^\mtc{A}}$.
        $$
            +,\ -,\ \cdot \in \mtc{F}_2 \implies +^\mtc{A},\ -^\mtc{A},\ \cdot^\mtc{A}: A^2 \to A
        $$
        $$
            \leq \in R_2 \implies \leq^\mtc{A} \in A^2;\ \ \ \ \ 0, 1 \in \mtc{C} \implies 0^\mtc{A}, 1^{\mtc{A}} \in A
        $$
    \end{itemize}
\end{eg}

\begin{obs}
    Podríamos intentar interpretar el lenguaje de los reales del ejemplo \ref{eg:lenguajes} con el universo $A = \C$, sin embargo, aunque podemos interpretar $+, -, \cdot, 0$ y $1$ de la forma habitual, no existe una interpretación de $\leq$ en $\C$.
\end{obs}

\begin{eg}[Lenguajes comunes]
    \begin{itemize}
        \item $L_\varnothing = \{\}$. Es el \textit{lenguaje vacío}, sigue teniendo símbolos generales. Sirve para expresar propiedades tales como: \textit{Existen tres elementos} ($\exists x_1, x_2, x_3$).
        \item $L_< = \{<\}$. Lenguaje para \textit{conjuntos ordenados}. Con $<\ \in R_2$.
        \item $L_{grupos} = \{+, -, 0\}$ (aditivo), $\{\cdot, ^{-1}, 1\}$ (multiplicativo). Lenguaje para \textit{grupos}. Con $+,\cdot\ \in \mtc{F}_2$; $-, ^{-1}\ \in \mtc{F}_1$; $0, 1 \in \mtc{C}$.
        \item $L_{cuerpos} = \{+, -, \cdot, 0, 1\}$. Lenguaje para \textit{cuerpos}.
        \item $L_{aritm\'etica} = \{+, \cdot, 0, 1, \leq\}$. Lenguaje para la \textit{aritmética}. También podemos considerar añadir otro símbolo de función $S$, cuya interpretación natural sería la función \textit{sucesor}.
        \item $L_{conj} = \{\in\}$. Lenguaje para \textit{conjuntos}. Todo se puede escribir con este lenguaje.
    \end{itemize}
\end{eg}

\subsection{Subestructuras}
\begin{dfn}[Subestructura de una L-estructura]
    Sean $\mtc{B}, \mtc{A}$ L-estructuras (con universos $B$ y $A$ respectivamente), decimos que $\mtc{A}$ es una subestructura de $\mtc{B}$ ($\mtc{A} \subseteq \mtc{B}$) si:
    \begin{itemize}
        \item $f^\mtc{A} = \left. f^\mtc{B} \right|_{A^n}$ para cada $f \in \mtc{F}_n$.
        \item $R^\mtc{A} = R^\mtc{B} \cap A^m$ para cada $R \in R_m$.
        \item $C^\mtc{A} = C^\mtc{B}$
        \item $A \subseteq B$
    \end{itemize}
\end{dfn}
\begin{eg}[Subestructuras. Ejemplos]
    Sean los lenguajes: $L_1 = \{+, 0\}$, $L_2 = \{+, -, 0\}$. Vamos a considerar las $L_1$-estructuras:
    $$
        W = \gen{\N, +^W, 0^W};\ Z = \gen{\Z, +^Z, 0^Z}
    $$
    Donde es fácil ver que se cumplen las condiciones de subestructura y podemos afirmar que $W \subseteq Z$. Sin embargo, si consideramos las $L_2$-estructuras:
    $$
        W' = \gen{\N, +^{W'}, -^{W'}, 0};\ Z' = \gen{\Z, +^{Z'}, -^{Z'}, 0}
    $$
    Vamos a dar una definición de $-^{W'}$ ya que el opuesto no está bien definido en $\N$. $-^{W'}: n \to 0$. Con esta interpretación es fácil ver que no se cumple que $-^{W'} = \left. f^{Z'} \right|_{\N}$ ya que:
    $$
        -^{W'}(2) = 0 \text{ y sin embargo } -^{Z'}(2)=-2
    $$
\end{eg}
\begin{obs}
    Consideremos el lenguaje $L_< = \{<\}$ de conjuntos ordenados, y las L-estructuras $\mtc{A} = \gen{A, <}$, $\mtc{B} = \gen{B, <}$ (que sólo tienen una relación). Es fácil ver que:
    $$
        \mtc{A} \subseteq \mtc{B} \iff A \subseteq B \text{ (suponemos $<^\mtc{A} = <^\mtc{B} \cap\ A^2$)}
    $$
\end{obs}

\section{Homomorfismos}
\begin{dfn}[Homomorfismo y monomorfismo]
    Sean $\mtc{A}$, $\mtc{B}$ L-estructuras y $\varphi: A \to B$ una función, decimos que $\varphi: \mtc{A} \to \mtc{B}$ es un \textbf{homomorfismo} de $\mtc{A}$ en $\mtc{B}$ si:
    \begin{enumerate}
        \item $\forall a_1, \ldots, a_n \in A$ se cumple: $\varphi(f^\mtc{A}(a_1, \ldots, a_n)) = f^\mtc{B}(\varphi(a_1), \ldots, \varphi(a_n))$
        \item $\forall a_1, \ldots, a_n \in A$ se cumple: $(a_1, \ldots, a_m) \in R^\mtc{A} \implies (\varphi(a_1), \ldots, \varphi(a_m) \in \mtc{B})$
        \item $\forall c \in \mtc{C}$ se cumple: $\varphi(c^\mtc{A}) = c^\mtc{B}$
    \end{enumerate}
    Además, si la segunda condición resulta ser necesaria y suficiente, $\varphi$ es un \textbf{monomorfismo}.
\end{dfn}
\begin{eg}[Ejemplo de monomorfismo]
    Sean $\mtc{A} = \gen{A, +, 0}$, $\mtc{B} = \gen{B, +, 0}$ y $\varphi(a+b) = \varphi(a) + \varphi(b);\ \varphi(0)=0$. Es fácil ver que $\varphi$ es un homomorfismo. Además, la única relación es la igualdad (que existe aunque no se especifique) y como $\varphi(a) = \varphi(b) \iff a = b$ estamos ante un monomorfismo.
\end{eg}

% !TeX root = ../logica-matematica.tex

\chapter{Términos y fórmulas de un lenguaje}

\section{Introducción}

Con los elementos de un lenguaje podemos definir sucesiones finitas de símbolos del lenguaje. Pueden ser:
\begin{itemize}
    \item \textbf{Términos}. Que pueden ser \textit{fuonciones} o \textit{elementos}.
    \item \textbf{Fórmulas}. Que pueden ser \textit{propiedades} o \textit{subconjuntos} de un universo.
\end{itemize}

\section{Términos}
\begin{dfn}[Términos de un lenguaje]
    Sea $L$ un lenguaje, y $M(L)$ el conjunto de palabras (sucesiones finitas de elementos de $L$). El conjunto de \textbf{términos} de $L$, $\ter(L)$, es el menor subconjunto de $M(L)$ que contiene a las constantes, variables y es cerrado para cada función $n$-aria $f \in \mtc F_n$ con $n \in \N$, es decir:
    $$
        \text{si } t_1, \ldots, t_n \in \ter(L) \implies f t_0 t_1 \ldots t_n \in \ter(L)
    $$
\end{dfn}

\begin{eg}[Ejemplos de términos]$ $
    \begin{itemize}
        \item $L = \sdf{+, -, \cdot, 0, 1}$.\\
        Son términos:
        \begin{itemize}
            \item $+ v_0 v_1$
            \item $0$
            \item $+ + v_0 v_1 + 1\ v_2$
        \end{itemize}
        \item $L = \sdf{\cdot, 1}$.\\
        Son términos:
        \begin{itemize}
            \item $v_0$
            \item $1$
            \item $\cdot v_0 v_1$
            \item $\cdot \cdot v_0 v_1 1$
        \end{itemize}
    \end{itemize}
\end{eg}

\begin{obs}
    Si consideramos el lenguaje $L = \sdf{+, -, 0, 1}$, vemos que también pertenecen al conjunto de términos los siguientes:
    $$
        0, 1, + 1\ 1, + + 1\ 1\ 1, + + 1\ 1 + 1\ 1, \ldots
    $$
    Y por tanto vemos que los enteros son \textit{términos abreviados} del lenguaje. ($2 \equiv + 1\ 1$).
\end{obs}

\begin{obs}[Notación]
    Durante el curso usaremos:
    \begin{itemize}
        \item $x, y, z$ para referirnos a variables $\sdf{v_i}_{i\in \N}$.
        \item $a, b, c$ para referirnos a elementos del universo de una estructura.
        \item $t(x_1, \ldots, x_n)$ para referirnos a un término $t \in \ter(L)$ donde aparecen algunas variables $x_1, \ldots, x_n$. Esto es una abreviatura, formalmente $t \in M(L)$ pero $t(x_1, \ldots, x_n) \notin M(L)$
    \end{itemize}
\end{obs}

\begin{dfn}[Función asociada a un término]
    Sea $t \in \ter(L)$, y $\mtc A$ una $L$-estructura, la \textbf{función asociada} al término $t$ es $t^\mtc A$, donde distinguimos ciertos casos:
    \begin{enumerate}
        \item Si $t$ es $t(x_1, \ldots, x_n)$, entonces $t^\mtc A: A^n \to A;\ (a_1,\ldots, a_n) \mapsto t^\mtc A (a_1, \ldots, a_n)$.
        \item Si $t$ es $x_i$, entonces $t^\mtc A(a_1, \ldots, a_n) = a_i$.
        \item Si $t$ es $c$, entonces $t^\mtc A(a_1, \ldots, a_n) = c^\mtc A$.
        \item Si $t$ es $f t_1 \ldots t_m$, entonces $t^\mtc A(a_1, \ldots, a_n) = f^\mtc A(t_1^\mtc A(a_1, \ldots, a_n), \ldots, t_m^\mtc A(a_1, \ldots, a_n))$.
    \end{enumerate}
\end{dfn}

\begin{obs}
    Se $t \in \ter(L)$ no tiene variables, entonces $t^\mtc A: A^n \to A$ es una función constante para cualquier $n$, es decir:
    $$
        t = t(x_1, x_2) = \cdots = t(x_1, \ldots, x_n)
    $$
\end{obs}

\begin{eg}[Ejemplo de funciones asociadas a términos]
    \begin{itemize}
        \item Sea $t: 1 + (1 + 1) \in \ter(L)$, y $\mtc Q = \gen{\Q, +, -, \cdot, 0, 1}$, si consideramos $t$ como $t(x_1, \ldots, x_n)$, entonces:
        $$
            t^\mtc Q: \Q^3 \to \Q;\ (a, b, c) \mapsto 3.
        $$
        \item Sea $t: 1 + x \in \ter(L)$, y $\mtc Q = \gen{\Q, +, -, \cdot, 0, 1}$, si consideramos $t$ como $t(x)$, entonces:
        $$
            t^\mtc Q: \Q \to \Q;\ 2 \mapsto 3;\ a \mapsto 1 + a
        $$
        Y además, si consideramos que $t$ es $t(x, y)$, entonces:
        $$
            t^\mtc Q: \Q^2 \to \Q;\ (2, 5) \mapsto 3;\ (a, b) \mapsto 1 + a
        $$
        \item Sea el lenguaje $L_\R = \sdf{+, -, \cdot, 0, 1} \cup \sdf{c_r}_{r\in \R}$, y sea $t: (x \cdot x) + (c_2 \cdot y) + (x \cdot y)$ (recordamos que $c_2$ se interpreta en $\mtc R_\R$ como $2 \in \R$), entonces la función asociada a $t(x, y)$ en la estructura $\mtc R_\R$ es:
        $$
            t^{\mtc R_\R}: \R^2 \to \R;\ (a, b) \mapsto (a \cdot a) + (c_2^{\mtc R_\R} \cdot b) + (a \cdot b) = a^2 + 2b + ab
        $$
    \end{itemize}
\end{eg}

\subsection{Sustitución}

\begin{dfn}[Término sustituido]
Sea el término $t: t(x_1, \ldots, x_n) \in \ter(L)$. Y sean $m$ términos en los que pueden aparecer las variables $x_1, \ldots, x_n$, $s_1: s_1(x_1, \ldots, x_n), \ldots, s_m: s_m(x_1, \ldots, x_n) \in \ter(L)$ con $m \leq n$.\\\\
Denotamos por $t(\quo{s_1}{x_1}, \ldots, \quo{s_m}{x_m})$ al \textbf{término sustituido}.
\end{dfn}

\begin{eg}[Sustitución en un término]
    Sean los términos $t: x + y$, $s: x + 1$, el término sustituido $t(\quo{s}{x})$ es:
    $$
        t(\quo{s}{x}) = (x + 1) + y
    $$
\end{eg}

\begin{pro}[Sustitución]
    Sean $x, y_1, \ldots, y_m \in V$ variables distintas, y $t: t(x, y_1, \ldots, y_m)$, $s: s(x, y_1, \ldots, y_m) \in \ter(L)$ términos, entonces $\forall (a, \mbf b) \in A^{m+1}$ se tiene que:
    $$
        \text{ la función asociada a $t(\quo{s}{x})$, } (t(\quo{s}{x}))^{\mtc A} \text{ evaluada en } (a, \mbf b) = t^\mtc A (s^\mtc A (a, \mbf b), \mbf b)
    $$
\end{pro}
\begin{proof}
    Para demostrarlo vamos a ver las equivalencias en una tabla:
    \begin{center}
        \begin{table}[ht]
            \centering
            \begin{tabular}{||c||c||c||c||}
                \hline
                $t$                & $t^\mtc A(s^\mtc A(a, \mbf b), \mbf b)$ & $t(s/x)$     & $t(s/x)^\mtc A (a, \mbf b)$ \\ \hline \hline
                $x$                & $s^\mtc A (a, \mbf b)$                  & $s$          & $s^\mtc A (a, \mbf b)$      \\
                $y_i$              & $b_i$                                   & $y_i$        & $b_i$                       \\
                $c$                & $c^\mtc A$                              & $c$          & $c^\mtc A$                  \\
                $f t_1 \ldots t_n$ & (\star)                                 & (\star\star) & (\star\star\star)           \\ \hline
            \end{tabular}
        \end{table}
    \end{center}

    \begin{itemize}
        \item[(\star)]
        $$
            f^\mtc A (t_1^\mtc A(s^\mtc A(a, \mbf b), \mbf b)), \ldots, t_n^\mtc A(s^\mtc A(a, \mbf b), \mbf b))
        $$
        \item[(\star\star)]
        $$
            f^\mtc A (t_1^\mtc A(s^\mtc A(a, \mbf b), \mbf b)), \ldots, t_n^\mtc A(s^\mtc A(a, \mbf b), \mbf b))
        $$
        \item[(\star\star\star)]
        $$
            f^\mtc A (t_1^\mtc A(s^\mtc A(a, \mbf b), \mbf b)), \ldots, t_n^\mtc A(s^\mtc A(a, \mbf b), \mbf b))
        $$
    \end{itemize}
    y como podemos observar, en todos las filas coinciden la segunda y la cuarta columna, con lo que hemos demostrado que son equivalentes.
\end{proof}

\subsection{Estructura generada por un subgrupo}

Sea $\mtc G = \gen{G, \cdot, ^-1, 1}$ un grupo, $L = \sdf{\cdot, ^-1, 1}$ un lenguaje y $S \subseteq G$ un subgrupo. Consideramos los $s_1, \ldots, s_n$ elementos del subgrupo $S$, podemos expresar la estructura asociada como:
$$
    s_1^{m_1}, \ldots, s_n^{m_n} = (x_1^{m_1}, \ldots, x_n^{m_n})^{\mtc G} (s_1, \ldots, s_n)
$$
Donde definimos $x_i^{m_i}$ como: $x_i \cdot (x_i \cdot \cdots (x_i \cdot x_i)))\ldots))$ (el producto de $x_i$ $m_i$ veces).\\\\

\begin{obs}
    Va a ser util definir ciertas abreviaturas para la expresión de los términos:
    \begin{itemize}
        \item $^{-1} t_i$ se abrevia por $t_i^{-1}$.
        \item $\cdot t_1 t_2$ se abrevia por $t_1 \cdot t_2$.
        \item $x_i^0$ se abrevia por $1$.
    \end{itemize}

\end{obs}

Una vez establecidas las abreviaturas, podemos definir la estructura generada por un subgrupo.

\begin{dfn}[Estructura generada por un subgrupo]
    Sea $\mtc G$ un grupo (como el definido anteriormente), y $S \subseteq G$, la \textbf{estructura generada} por el \textbf{subgrupo} $S$ es:
    $$
        \gen{S} = \sdf{t^\mtc G(s_1, \ldots, s_n) \mid t \in \ter(L),\ n \in \N,\ s_i \in S}
    $$
\end{dfn}

De esta noción surge la definición de subestructura generada por un conjunto.

\begin{pro}[Subestructura generada por un conjunto]
    Sea $L$ un lenguaje, $\mtc A$ una $L$-estructura, y $X \subseteq A$ un subconjunto. Si $X \neq \varnothing$ o en $L$ hay al menos una constante, el conjunto:
    $$
        \gen X_{\mtc A} = \sdf{t^\mtc A (a_1,\ldots, a_n) \mid t\in \ter(L); a_1,\ldots,a_n \in X \text{ con } n\in\N}
    $$
    es la mínima subestructura de $\mtc A$ que contiene a $X$.
\end{pro}

\begin{proof}
    Completar.
\end{proof}

\section{F\'ormulas}

Sea $L$ un lenguaje, $L = \sdf{f, \ldots, R, \ldots, c, \ldots}$ además del conjunto de variables $V = \sdf{v_i}_{i \in \N}$ y conectores lógicos: $=, \neg, \then, \forall$. En la sección anterior caracterizamos el conjunto de \textit{términos} del lenguaje. Ahora vamos a caracterizar las \textit{fórmulas}.

\begin{dfn}[Fórmulas de un lenguaje]
    Sea $L$ un lenguaje. El conjunto de $L$-fórmulas, $\for(L)$, es el mínimo subconjunto de $M(L)$ (palabras de $L$) que contiene a \textit{fórmulas atómicas} y es cerrado por $\neg, \then$ y $\forall v_i$ para cada $i \in \N$.\\

    Las \textit{fórmulas atómicas} pueden ser de dos tipos:
    \begin{itemize}
        \item $= t_1 t_2$ con $t_1, t_2 \in \ter(L)$.
        \item $R t_1 \ldots t_m$, para todo $t_i \in \ter(L)$ y $R \in \mtc R_m$.
    \end{itemize}
\end{dfn}

\begin{obs}
    Sean $F, G \in \for(L)$, entonces como $\for(L)$ es cerrado para los operadores lógicos:
    $$
        \neg F, \then F G\ (F \then G),\ \forall v_i F \text{ también son fórmulas de L}
    $$
    El resto de conectores lógicos los usamos como abreviaturas de $\neg, \then, \forall$. Por ejemplo:
    \begin{itemize}
        \item $F \lor G$ como abreviatura de $\neg F \then G$.
        \item $\exists v_i$ como abreviatura de $\neg \forall v_i \neq F$.
    \end{itemize}
\end{obs}

\begin{eg}[Ejemplos de fórmulas]
    Sea $L = \sdf{+, -, \cdot, 0, 1, \leq}$  y consideramos los términos abreviados de $L$: $t_1: 3x_1+2$, $t_2: x_1-x_3$, $t_3: y^2-2x+1$.
    \begin{itemize}
        \item $3x_1 + 2 = x_1 - x_3$, es una fórmula atómica abreviada ($= t_1 t_2$).
        \item $3x_1 + 2 \leq x_1 - x_3$, es una fórmula atómica abreviada ($\leq t_1 t_2$).
        \item $\forall x \exists y\ x\leq y$, es una fórmula abreviada.
        \item $\exists y\ x \leq y$, es una fórmula abreviada ($\neg \forall v_0\ \neg \leq v_0 v_1$).
        \item $\exists y\ y^2-2x+1 = 0$.
    \end{itemize}
\end{eg}

\begin{dfn}[Subfórmula]
    Una \textbf{subfórmula} $G$ de un fórmula $F \in \for(L)$ es una fórmula que aparece en $F$.
\end{dfn}

\begin{eg}[Ejemplos de subfórmulas]
    Sea $F \in \for(L), F: \forall x\ \exists y\ x\leq y$, las subfórmulas que aparecen en $F$ son:
    $$
        \forall x\ \exists y\ x\leq y,\ \exists y\ x\leq y,\ x\leq y
    $$
\end{eg}

\begin{eg}[Árbol de decisión de una fórmula]
    Sea $F \in  \for(L)$, $F: (\exists x\ x\leq 1) \lor (\forall y\ y+1=3)$. Podemos representar su árbol de decisión, cuyos nodos son las subfórmulas de $F$.\\
    \begin{center}
        \begin{tikzpicture}
            %\tikzset{edge from parent/.style={draw,edge from parent path={(\tikzparentnode.south)-- +(0,-8pt)-| (\tikzchildnode)}}}
            \Tree [.$$(\exists x\ x\leq 1)\lor (\forall y\ y+1=3)$$
            [.$$(\exists x\ x\leq 1)$$
            [.$$\exists x$$ ]
            [.$$x\leq 1$$ ] ]
            [.$$(\forall y\ y+1=3)$$
            [.$$\forall y$$ ]
            [.$$y+1=3$$ ] ] ]
        \end{tikzpicture}
    \end{center}
\end{eg}

\subsection{Variables libres y ligadas}
Sea $F \in \for(L)$ una fórmula en la que aparecen algunas variables, éstas pueden aparecer asociadas, libres o ligadas. Antes de dar una definición veamos un ejemplo.

\begin{eg}[Ejemplos de tipos de variables]
    \begin{itemize}
        \item La fórmula $\forall x\ 0\leq x$ expresa \textit{$0$ es mínimo}. En esta fórmula, $x$ aparece dos veces, la primera está \textit{asociada} al operador, y la segunda es una \textit{variable ligada}
        \item Consideramos la fórmula: $(\forall x\ 0\leq x) \lor x = 3$. En ella $x$ aparece tres veces, la primera está \textit{asociada} al operador, y la segunda es una \textit{variable ligada} y la última es una \textit{aparición libre}.
    \end{itemize}
\end{eg}

\begin{dfn}[Aparición de una variable. Tipos]
    Sea $L$ un lenguaje, $F \in \for(L)$ una fórmula y $x\in V$ una variable.
    \begin{enumerate}
        \item Decimos que $x$ \textbf{aparece} en $F$ si $x$ es algunos de los símbolos de $F$.
        \item Una aparición de $x$ en $F$ está \textbf{asociada} a $\forall$ si la aparición va precedida de $\forall$.
        \item Una aparición de $x$ en $F$ es \textbf{ligada} si $x$ aparece en una subfórmula $G$ de $F$ tal que $\forall x G$ es una subfórmula de $F$.
        \item Una aparición de $x$ en $F$ es \textbf{libre} si no está \textit{asociada} a $\forall$ y no es \textit{ligada}.
    \end{enumerate}
\end{dfn}

\begin{obs}[Notación]
    Sea $F \in \for(L)$, escribimos $F(x_1, \ldots, x_n)$ si las variables \textit{libres} de $F$ están en $\sdf{x_1, \ldots, x_n}$. Por ejemplo:
    \begin{itemize}
        \item $F(x): \exists y\ x = 2y$, expresa \textit{x es par}.
        \item Podemos incluso sustituir variables libres por otros términos, si $s$ es $x+1$ entonces:
        $$
            F(\quo{s}{x}): \exists y\ x+1=2y, \text{ que expresa \textit{x+1 es par}}
        $$
        Sin embargo, hay sustituciones que aunque sintácticamente podemos hacerlas, no tienen sentido semánticamente, por ejemplo:
        $$
            F(\quo{y}{x}): \exists y\ y=2y, \text{ que no expresa \textit{y es par}}
        $$
    \end{itemize}
\end{obs}

\begin{dfn}[Variable sustituible]
    Sea $L$ un lenguaje, $F \in \for(L)$, $x \in V$, $s \in \ter(L)$.
    \begin{enumerate}
        \item $x$ es \textbf{sustituible} por $s$ en $F$ si no hay una aparición libre de $x$ en una subfórmula de $F$ de la forma $\forall y G$ donde $y$ es una variable de $s$.
        \item Si $x$ es \textit{sustituible} por $s$ en $F$, $F(\quo{s}{x})$ o $F[\quo{s}{x}]$ denota la fórmula obtenida de $F$ sustiyendo \textit{simultanemante} las apariciones \textit{libres} de $x$ por $s$.
    \end{enumerate}
\end{dfn}

\begin{eg}[Ejemplos de sustitución en fórmulas]$ $
    \begin{itemize}
        \item Sea $F: \forall y\ x = z$, $x \in V$ y $s: x+1 \in \ter(L)$.\\
        $$
            F[\quo{s}{x}]: \forall y\ x+1 = z
        $$
        \item Sea $F: (\exists x\ x=z) \lor (x+y=1)$, $x \in V$ y $s: y+1 \in \ter(L)$.\\
        $$
            F[\quo{y+1}{x}]: (\exists x\ x=z) \lor (y+1+y=1)
        $$
    \end{itemize}
\end{eg}

\begin{obs}
    Si $s \in \ter(L)$ no tiene variables, entonces se puede sustituir en cualquier variable en cualquier fórmula.
\end{obs}

\begin{dfn}[Cierre universal]
    Sea $F \in \for(L)$, un \textbf{cierre universal} de $F$ es una fórmula obtenida precediendo a $F$ con $\forall x_1 \ldots \forall x_n$ si $F$ es $F(x_1, \ldots, x_n)$.
\end{dfn}

\begin{eg}[Ejemplo de cierre universal]
    Sea $F: x \leq 2$, cierres universales de $F$ pueden ser: $\forall x\ x\leq 2$, $\forall x \forall y \forall z\ x \leq 2$.
\end{eg}

\begin{dfn}[Enunciado]
    Una fórmula se llama \textbf{enunciado} (o fórmula cerrada) si no tiene variables libres.
\end{dfn}

\begin{obs}
    Los cierres universales de fórmulas son enunciados.
\end{obs}

\begin{dfn}[Fórmula existencial]
    Una fórmula abreviada es \textbf{existencial} si es de la forma $\exists x_1 \ldots \exists x_n F$ con $F$ sin cuantificadores.
\end{dfn}
\begin{eg}[Ejemplo de fórmula existencial]
    Una fórmula existencial puede tener variables libres. Por ejemplo:
    $$
        F(z):  \exists x_1 \exists x_2\ (x_1 = x_2 \lor x_1 + 3 = z)
    $$
\end{eg}

% !TeX root = ../logica-matematica.tex
\chapter{Relación de satisfacción.}

En esta sección vamos a hablar de la satisfacibilidad de fórmulas en distintas estructuras. Como ejemplo introductorio, sea $L = \sdf{\leq}$ un lenguaje, $\mtc Z = \gen{\Z, \leq}$, $\mtc N = \gen{\N, \leq}$ $L$-estructuras y $E: \forall x\ 0 \leq x$ una $L$-fórmula.\\
Vemos que se cumple (formalmente, se \textit{satisface}) en $\mtc N$ pero no en $\mtc Z$, ya que $-1 \in \Z$ no satisface $E$ en $\mtc Z$.

\section{Satisfacibilidad}

\begin{dfn}[Satisfacibilidad de una fórmula]
    Sea $L$ un lenguaje, $\mtc A$ una $L$-estructura, $F(x_1, \ldots, x_n) \in \for(L)$, $\mbf a = (a_1, \ldots, a_n) \in A^n$ y $\mbf x = (x_1, \ldots, x_n)$. Decimos que $a$ \textbf{satisface} $F$ en $\mtc A$ ($\mtc A \vDash F(\mbf a)$), o $\mtc A$ satisface $F$ cuando a la variable $\mbf x$ se le asigna el valor $\mbf a$, o $F(\mbf a)$ se satisface en $\mtc A$ si:

    \begin{itemize}
        \item Si $F$ es $R t_1 \ldots t_l$ y $(t_1^\mtc A(\mbf a), \ldots, t_l^\mtc A(\mbf a)) \in \mtc R^\mtc A$.
        \item Si $F$ es $\neg G$ y $\mtc A \nvDash G(\mbf a)$.
        \item Si $F$ es $G \then H$ y si $\mtc A \vDash G(\mbf a)$ entonces $\mtc A \vDash H(\mbf a)$.
        \item Si $F$ es $\forall v G$ con $v \neq x_i$, $G(v, x_2, \ldots, x_n)$ y $\forall d \in A$ se tiene que: $\mtc A \vDash G(d, \mbf a)$, es decir $\forall d \in A$, $\mtc A \vDash G(d, a_2, \ldots, a_n)$.
    \end{itemize}
\end{dfn}

\begin{eg}[Ejemplo de relación de satisfacción]
    Sea $L = \sdf{+, -, \cdot, 0, 1, <}$ y $\mtc R = \gen{\R, +, -, \cdot, 0, 1, <}$. Sea $F: \exists x (x < 2 \land 4 < x^2)$ una fórmula abreviada, ¿existe $(a, b) \in \R^2$ tales que $\mtc R \vDash F(a, b)$?\\

    $F$ sin abreviar sería:
    $$
        F: \neg \forall x\ (x < 2 \then \neg 4 < x^2)
    $$
    Entonces tenemos:
    \begin{align*}
        \mtc R \vDash F(a, b) &\iff \mtc R \nvDash \forall x (x < 2 \then \neg 4 < x^2)(a, b)\\
                               &\iff\exists d \in \R,\  \mtc R \nvDash \forall x (x < 2 \then \neg 4 < x^2)(d, b)\\
                               &\iff \exists d \in \R,\ \mtc R \vDash (x < 2)(d, b) \text{ y } \mtc R \nvDash(\neg 4 < x^2)(d, b)\\
                               &\iff \exists d \in \R,\ d < 2 \text{ y } \mtc R \vDash (4 < x^2)(d, b)\\
                               &\iff \exists d \in \R,\ d < 2 \text{ y } 4 < d^2; \text{ se satisface tomando $d = -3 \in \R$}
    \end{align*}
    Con lo que concluimos que: $\mtc R \vDash F(-3, b)$.
\end{eg}

\begin{pro}[$\vDash$ solo depende de las variables libres]\label{pro:1.1}
    Sea $L$ un lenguaje, $\mbf x = (x_1, \ldots, x_n)$ e $\mbf y = (y_1, \ldots, y_m)$, con $x_i, y_i \in V$ $x_i \neq y_i$. Sean $F \in \for(L)$, $F(\mbf x, \mbf y)$ y $\mtc A$ una $L$-estructura; entonces para todo $\mbf a = (a_1, \ldots, a_n) \in A^n$:
    $$
        \mtc A \vDash F(\mbf a) \iff \text{ existe } \mbf b \in A^m : \mtc A \vDash F(\mbf a, \mbf b) \iff \forall \mbf b \in A^m \mtc A \vDash F(\mbf a, \mbf b)
    $$
\end{pro}

\begin{proof}
    A completar. %% TODO: Página: 13, día 26/09/2019
\end{proof}

\begin{dfn}[Satisfacibilidad de un enunciado]
    Sea $E \in \for(L)$ un enunciado (sin variables libres), $\mtc A$ una $L$-estructura. Decimos que $\mtc A$ \textbf{satisface} $E$, o $E$ es verdadero en $\mtc A$, o $E$ se satisface en $\mtc A$ si al considerar $E$ como $E(y)$ se tiene que:
    $$
        \mtc A \vDash E(b) \text{ para todo (algún) } b\in A
    $$
    Lo escribimos como $\mtc A \vDash E$
\end{dfn}

\begin{obs}
    En la definición hablamos de que el enunciado se satisfaga para algún $b \in A$ o todo $b \in A$. Esto es equivalente ya que siendo $E$ un enunciado:
    $$
        \exists b \in A,\ \mtc A \vDash E(b) \iff \forall b \in A,\ \mtc A \vDash E(b) \text{ por la proposición } \ref{pro:1.1}
    $$
\end{obs}

\subsection{Validez de fórmulas}

\begin{dfn}[Fórmula válida]
    Sea $L$ un lenguaje:
    \begin{enumerate}
        \item Sea $F$ un enunciado de $L$, decimos que $F$ es una \textbf{fórmula válida} si $\mtc A \vDash F$ para toda $L$-estructura $\mtc A$.
        \item Sea $G \in \for(L)$, $G$ es una fórmula válida si $G$ es $G(x_1, \ldots, x_n)$ y $\forall x_1 \ldots x_n G$ es un enunciado válido.
    \end{enumerate}
\end{dfn}

\section{Conjuntos definibles}

\begin{dfn}[Conjunto definible]
    Sea $L$ un lenguaje, $\mtc A$ una $L$-estructura:
    \begin{itemize}
        \item Un conjunto $X \subseteq A^n$ para algún $n$ es \textbf{definible} (en $\mtc A$) si existe una fórmula $F(x_1, \ldots x_n) \in \for(L)$ tal que:
        $$
            X = \sdf{\mbf a \in A^n \mid \mtc A \vDash F(\mbf a)}
        $$
        \item Un conjunto $X \subseteq A^n$ para algún $n$ es \textbf{definible con parámetros} si existe una fórmula $F(x_1, \ldots x_n) \in \for(L_A)$ tal que:
        $$
            X = \sdf{\mbf a \in A^n \mid \mtc A_A \vDash F(\mbf a)}
        $$
    \end{itemize}
\end{dfn}

\begin{eg}[Ejemplos de conjuntos definibles]$ $
    Sea $L = \sdf{+, -, \cdot, 0, 1}$, $\mtc C = \gen{\C, +, -, \cdot, 0, 1}$. Sea $f(x_1, \ldots, x_n) \in \C[x_1, \ldots, x_n]$, entonces:
    \begin{enumerate}
        \item
        $$
            X = \sdf{\mbf a \in \C^n \mid f(\mbf a) = 0} \text{ es un conjunto definible con parámetros}
        $$
        La fórmula $F$ sería, $F(x_1, \ldots, x_n): f(x_1, \ldots, x_n) = 0$ con $F \in \for(L_\C)$.

        \item $$X = \sdf{(a, b) \in \C^2 \mid a\cdot b = 1} \text{ es definible en $\mtc C$ sin parámetros}$$

        \item ¿Es $\sdf{-\sqrt 2, \sqrt 2}$ definible en $\mtc \C$? Sí, la fórmula (abreviada) que lo define es $x^2 = 2$. ($ = \cdot\ x\ x\ +\ 1\ 1$).
    \end{enumerate}
\end{eg}

\begin{obs}
    Si permitimos parámetros, cualquier subconjunto finito de $\C^n$ es definible. Sea el conjunto:
    $$
        X = \sdf{(a_{11}, \ldots, a_{1n}), \ldots, (a_{m1}, \ldots, a_{mn})} \text{ con } a_{ij} \in \C
    $$
    En el caso $m = 1$, basta considerar la fórmula $F: x_1 = a_{11} \land x_2 = a_{12} \land \cdots \land x_n = a_{1n}$.\\
    En el caso general, $F_i(x_1, \ldots, x_n): x_1 = a_{i1} \land x_2 = a_{i2} \land \cdots \land x_n = a_{in} \in \for(L_a)$. Con esto definimos $X$ como:
    $$
        X = \sdf{\mbf b = (b_1, \ldots, b_m) \in \C^n \mid \mtc C \vDash (F_1 \lor \cdots F_m)(\mbf b)}
    $$
\end{obs}

\begin{eg}
    Sea $L = \sdf{S, +, \cdot, 0, <}$ y $\mtc N = \gen{\N, S, +, \cdot, 0, <}$. Podemos definir el conjunto de números pares naturales. Consideramos la $L$-fórmula:
    $$
        \text{ $x$ es par }: \exists y\ (x = y + y)
    $$
    entonces, definimos el conjunto $X$ de naturales pares como:
    $$
        X = \sdf{ m \in \N \mid \N \vDash (\exists y\ (x = y + y))(m)}
    $$
\end{eg}

\begin{obs}
    Sea $\mtc R$ el cuerpo de los reales. Sea $F: y = 0$, $G: x = 0$, con $F(x, y)$ y $G(x, y)$. Se consideran los conjuntos:
    \begin{align*}
        X_1 &= \sdf{(a, b) \in \R^2 \mid \mtc R \vDash F(a, b)}\\
        X_2 &= \sdf{(a, b) \in \R^2 \mid \mtc R \vDash G(a, b)}
    \end{align*}
    Entonces $X_1 \bigcup X_2$ es definible mediante $F \lor G$.
\end{obs}

\begin{pro}[Caracterización de las operaciones de conjuntos]
    Sea $L$ un lenguaje, $\mtc A$ una $L$-estructura, $X, Y \subseteq A^n$ definibles mediante $F(x_1, \ldots, x_n)$ y $G(x_1, \ldots, x_m)$ respectivamente. Entonces:
    \begin{itemize}
        \item $ X \bigcup Y$ es definible mediante $F \lor G$.
        \item $ X \bigcap Y$ es definible mediante $F \land G$.
        \item $A^n \setminus X$ es definible mediante $\neg F$.
        \item Sea $Z$ la proyección de $X$ en las primeras $n - 1$ coordenadas:
        $$
            Z = \sdf{(a_1, \ldots, a_{n-1}) \in A^{n-1} \mid \text{ existe } b\in A:\ (a_1, \ldots, a_{n-1}, b) \in X}
        $$
        es decir,
        $$
            Z = \sdf{(a_1, \ldots, a_{n-1}) \in A^{n-1} \mid \mtc A \vDash (\exists x_n F)(a_1, \ldots, a_{n-1})}
        $$
        \item Sea $W \subseteq A^m$ definible mediante $H(x_1, \ldots, x_m)$.
        $$
            X \times W \subseteq A^{n+m} \text{ es definible}
        $$
        Vamos a necesitar una fórmula de la forma $F \land H$, sin embargo, eso definiría la intersección. Notamos que $H': H(\quo{y_1}{x_1}, \ldots, \quo{y_m}{x_m})$ también define $W$. Con ello obtenemos que:
        $$
            X \times W = \sdf{(\mbf a, \mbf b) \in A^{n+m} \mid \mtc A \vDash (F \land H')(\mbf a, \mbf b)}
        $$
    \end{itemize}
\end{pro}
\begin{proof}
    Son propiedades que no se han demostrado en clase.
\end{proof}

\begin{pro}[Satisfacibilidad de una sustitución]\label{pro:1.2}
    Sea $L$ un lenguaje, $s \in \ter(L)$ con $s$ siendo $s(x, y_1, \ldots, y_m)$ e $y_i$ que aparece en $s$. Sea $F \in \for(L)$, $F$ es $F(x, y_1, \ldots, y_m, z_1, \ldots, z_l)$ con $(z_j \neq y_i)$, tal que $x$ es sustituible por $s$ en $F$. (Es decir, no hay subfórmulas de $F$ de la forma $\forall y_i G$ con $x$ libre en $F$).\\\\
    Entonces para toda $L$-estructura $\mtc A$ y $(a, \mbf b, \mbf c) \in A^{1+m+l}$ se tiene que:
    $$
        \mtc A \vDash F(\quo{s}{x})(a, \mbf b, \mbf c) \iff \mtc A \vDash F(s^\mtc A(a, \mbf b), \mbf b, \mbf c)
    $$
\end{pro}

\begin{proof}
    A completar. %%TODO: Página 16, día 30/09/2019
\end{proof}

\begin{cor}[a la proposición \ref{pro:1.2}. Cambio de nombre a una variable ligada.]
    Sea $L$ un lenguaje, $F(x, \mbf z) \in \for(L)$, $y \in V$ que no aparece en $F$. Entonces para toda $L$-estructura $\mtc A$ y $\mbf \in A^m$:
    $$
        \mtc A \vDash \forall y (F(\quo{y}{x}))(\mbf b) \iff \mtc A \vDash (\forall x F)(\mbf b)
    $$
    En particular si $F$ es $F(x)$:
    $$
        \forall y\ F(\quo{y}{x}) \bicond \forall x\ F \text{, es válida.}
    $$
\end{cor}
\begin{proof}
    A completar. %%TODO: Página 17, día 01/10/2019
\end{proof}

% !TeX root = ../logica-matematica.tex

\chapter{Tautologías}

Recordemos los valores de verdad de algunas construcciones con $P, Q$ variables proposicionales.\\
\begin{center}
    \begin{table}[ht]
    \centering
    \begin{tabular}{c|c|c|c|c}
    $P$ & $Q$ & $P \lor Q$ & $P \then Q$ & $P \then P$ \\ \hline
    1   & 1   & 1          & 1           & 1           \\
    1   & 0   & 1          & 0           & 1           \\
    0   & 1   & 1          & 1           & 1           \\
    0   & 0   & 0          & 1           & 1
    \end{tabular}
    \end{table}
\end{center}

El hecho de que en las implicaciones sólo haya un valor falso nos será de utilidad para construir \textit{tautologías}.

\begin{dfn}[Formulas básica]
    Sea $L$ un lenguaje, una \textbf{formula básica} de $L$ es una formula o \textit{atómica} o de la forma: $\forall x G$ con $G \in \for(L)$.
\end{dfn}

\begin{dfn}[Subfórmulas básicas necesarias]
    Sea $L$ un lenguaje, $F \in \for(L)$, el conjunto de \textbf{subfórmulas básicas necesarias} de $F$, $\sbn(F)$ es:
    \begin{itemize}
        \item Si $F$ es básica: $\sbn(F) = \sdf{F}$.
        \item Si $F$ es $\neg G$: $\sbn(F) = \sbn(G)$.
        \item Si $F$ es $G \then H$: $\sbn(F) = \sbn(G) \bigcup \sbn(H)$.
    \end{itemize}
\end{dfn}

\begin{obs}
    Toda fórmula se obtiene de subfórmulas básicas necesarias usando $\neg$ y $\then$.
\end{obs}

\begin{dfn}[Distribución de valores de verdad]
    Sea $L$ un lenguaje.
    Una \textbf{distribución de valores de verdad} (d.v.v.) de $L$ es una aplicación:
    $$
        \delta: \sdf{F \in \for(L) \mid F \text{ es básica}} \to \sdf{0, 1}
    $$
\end{dfn}

\begin{dfn}[Tautología]
    Una \textbf{tautología} de $L$ es una fórmula $F \in \for(L)$ tal que $\bar \delta(F) = 1, \forall \delta$ (con $\delta$ una d.v.v. ) donde $\bar \delta$ es la única aplicación $\bar\delta: F \to \sdf{0, 1}$ que extiende a $\delta$ y satisface:
    \begin{itemize}
        \item $\bar\delta(\neg G) = 1$ si $\bar\delta(G) = 0$ y $\bar\delta(\neg G) = 0$ si $\bar\delta(G) = 1$
        \item $\bar\delta(G \then H) = 0$ si $\bar\delta(G) = 0 \land \bar\delta(H) = 1$ y $\bar\delta(G \then H) = 1$ en el resto de casos.
    \end{itemize}
\end{dfn}

\begin{dfn}[Fórmulas tautológicamente equivalentes]
    Sea $L$ un lenguaje, $F, G \in \for(L)$ fórmulas de $L$. Decimos que $F$ y $G$ son \textbf{tautológicamente equivalentes} si:
    $$
        F \bicond G \text{ es una tautología.}
    $$
\end{dfn}

\begin{obs}
    Para fórmulas abreviadas tomamos $\exists x G$ como fórmula básica también.
\end{obs}

\begin{eg}[Comprobación de dos fórmulas tautológicamente equivalentes]
    Sean las fórmulas $G: F_1 \then (F_2 \then F_3)$ y $H: (F_1 \land F_2) \then F_3$. Vamos a ver que $G$ y $H$ son tautológicamente equivalentes.\\\\
    Tenemos que comprobar que $G \bicond H$ es una tautología. Es decir, $G \then H$ y $H \then G$ son tautologías. Supondremos que no lo son y al asignar los valores de verdad correspondientes llegaremos a una contradicción.\\
    $$
    \begin{matrix}
        F_1 & \then & (F_2 & \then & F_3) & \then & (F_1 & \land & F_2) & \then & F_3 \\
            & \mbf 1&      &       &      &\mbf 0 &      &       &      &\mbf 0 &     \\
            &   1   &      &       &      &   0   &      &\mbf 1   &      &    0  &\mbf 0  \\
            &   1   &      &       &      &   0   &\mbf 1   &   1   &\mbf 1   &    0  &  0  \\
    \mbf 1   &   1   &\mbf 1  &       &\mbf 0   &   0   &  1   &   1   &  1   &    0  &  0  \\
         1   & \mbf 0&   1  &\mbf 0   &  0   &   0   &  1   &   1   &  1   &    0  &  0  \\
    \end{matrix}
    $$
    Se marcan en negrita las deducciones en cada paso (fila). Podemos ver que llegamos a una contradicción en el valor de verdad de la primera implicación, con lo que la expresión es una tautología. Se procede de forma análoga para el recíproco.
\end{eg}

\begin{pro}[Tautologías y fórmulas válidas]\label{pro:1.3}
    Sea $L$ un lenguaje. Las tautologías de $L$ son fórmulas válidas de $L$.
\end{pro}
\begin{proof}
    A completar. %% TODO: Página 18, día 01/10/2019
\end{proof}

\begin{obs}
    La propiedad de ser tautología es puramente sintáctica.\\
    $$
        \forall x F \then F \text{ no es una tautología.}
    $$
\end{obs}

% !TeX root = ../logica-matematica.tex

\chapter{Consecuencia semántica}

\begin{dfn}[$L$-teoría]
    Sea $L$ un lenguaje. Una \textbf{$L$-teoría} es un conjunto de \textit{enunciados} de $L$ (no tiene por qué ser finito).
\end{dfn}

\begin{eg}[Ejemplos de $L$-teorías]
    \begin{enumerate}
        \item Sea $L = \sdf{\cdot, 1}$ el lenguaje de grupos. Sean $G_1, G_2, G_3$ los axiomas de grupos. $T = \sdf{G_1, G_2, G_3}$ es una $L$-teoría.
        \item GADST: (grupos abelianos divisibles sin torsión). Sea $L = \sdf{+, 0}$ y las propiedades:
            \begin{itemize}
                \item $G_1: +$ es asociativa.
                \item $G_2: 0$ es el neutro.
                \item $G_3: \forall x \exists y\ x+y = 0$.
                \item $G_4: \forall x \forall y\ x+y=y+x$.
                \item $G_{5n}: \forall x \exists y\ x = n \cdot y$ ($n$-divisible). $n \cdot y$ abrevia $\sum_1^n y$.
                \item $G_{6n}: \forall\ x (n\cdot x = 0 \then x = 0)$
            \end{itemize}
            Entonces GADST $= \sdf{G_1, G_2, G_3, G_3} \bigcup \sdf{G_{5n}}_{n>1,\ n \in \N} \bigcup \sdf{G_{6n}}_{n>1,\ n \in \N}$.
        \item $\sdf{\forall x\ x \neq x}$ y $\sdf{F \land \neg F}$ son teorías, aunque sean \textit{incoherentes}.
    \end{enumerate}
\end{eg}

\begin{dfn}[Modelo]
    Sea $L$ un lenguaje, $T$ una $L$-teoría. Un \textbf{modelo} de $T$ es una $L$-estructura $\mtc A$ tal que:
    $$
        \mtc A \vDash F \text{ para toda } F \in T \text{ y se escribe } \mtc A \vDash T
    $$
\end{dfn}

\begin{eg}[Ejemplos de modelos]
    \begin{enumerate}
        \item $T_{\text{grupos}}$ en $L = \sdf{\cdot, 1}$, sea $\mtc G$ una $L$-estructura. $\mtc G$ es un modelo de $T_{\text{grupos}} \iff \mtc G$ es un grupo.
        \item $\mtc A \vDash$ GADST $\iff \mtc A$ es un grupo abeliano divisible sin torsión.
        \item $\sdf{\forall x\ x \neq x}$ y $\sdf{F \land \neg F}$ no tienen modelos.
    \end{enumerate}
\end{eg}

\begin{obs}[Notación]
    $\model(T)$ designa la clase de todos los modelos de $T$.
\end{obs}

\begin{obs}
    Sean $T_1, T_2$ teorías, si $T_1 \subseteq T_2$ entonces $\model(T_2) \subseteq \model(T_1)$.
\end{obs}

\begin{dfn}[Teoría de una estructura]
    Sea $\mtc A$ una $L$-estructura, una \textbf{teoría} de $\mtc A$ es el conjunto:
    $$
        \te(\mtc A) = \sdf{F \in \for(L) \mid F \text{ es un enunciado y } \mtc A \vDash F}
    $$
\end{dfn}

\begin{obs}$ $
    \begin{itemize}
        \item Sea $F$ un enunciado. $\mtc A \vDash F$ establece una relación semántica entre un objeto semántico $\mtc A$ y otro sintáctico $F$.
        \item Sea $T$ una $L$-teoría (objeto sintáctico), también escribimos $\mtc A \vDash T$.
    \end{itemize}
\end{obs}

\begin{dfn}[Consecuencia semántica]
    Sea $L$ un lenguaje, $T$ una $L$-teoría, y $F$ un $L$-enunciado. $F$ es \textbf{consecuencia semántica} de $T$ ($T \vDash F$) si $\mtc A \vDash F$ para toda estructura $\mtc A$ tal que $\mtc A \vDash T$.
\end{dfn}

\begin{obs}$ $
    \begin{itemize}
        \item Si $T = \varnothing$, escribiremos $\vDash F$.
        \item Sea $G \in \for(L)$, $G$ es consecuencia semántica de $T$ ($T \vDash G$) si $T \vDash F$ para algún (o todo) cierre universal $F$ de $G$.
        $$
            T \vDash \forall x_1 \ldots \forall x_n\ G \implies T \vDash G
        $$
    \end{itemize}
\end{obs}

%% TODO: Agregar ejemplos página 20, día 07/10/2019

\begin{dfn}[Fórmulas equivalentes respecto de una teoría]
    Sea $L$ un lenguaje, $T$ una $L$-teoría, $F_1, F_2 \in \for(L)$. Diremos que $F_1$ y $F_2$ son \textbf{equivalentes respecto de $T$} si:
    $$
        T \vDash \forall x_1 \ldots \forall x_n\ (F_1 \bicond F_2)
    $$
\end{dfn}

%% TODO: Agregar ejemplo página 21, día 07/10/2019

\begin{obs}
    Si $F_1, F_2 \in \for(L)$ son equivalentes con respecto a $T$, entonces $F_1$ y $F_2$ definen el mismo conjunto en cualquier modelo de $T$.
\end{obs}

\part{Sistemas formales de primer orden}
% !TeX root = ../logica-matematica.tex

\chapter{Consecuencia sintáctica}

Hemos visto que sea $F$ un enunciado de un lenguaje $L$, $\mtc A$ una $L$-estructura y $T$ una $L$-teoría, hemos definido relaciones semánticas, tanto entre objetos semánticos y sintácticos ($\mtc A \vDash F$) como entre objetos puramente sintácticos ($T \vDash F$).\\
En este bloque vamos a introducir la \textit{consecuencia sintática}, una relación sintáctica entre objetos sintácticos ($T \vdash F$).

\section{Sistema formal. Axiomas}
\begin{dfn}[Sistema formal]
    Nos referimos por \textbf{sistema formal} al conjunto $SF = \sdf{L, Ax, RDec, T}$ donde:
    \begin{itemize}
        \item $L$ es un lenguaje.
        \item $Ax$ es un conjunto de axiomas.
        \item $RDec$ es un conjunto de reglas de deducción que nos permiten obtener expresiones del lenguaje a partir de otras.
        \item $T$ es un conjunto de teoremas, expresiones del lenguaje que se deducen a partir de los axiomas utilizando las reglas de deducción.
    \end{itemize}
\end{dfn}

En nuestro curso, nuestro $SF$ consta de:
\begin{itemize}
    \item $L$ un lenguaje de primer orden.
    \item $Ax$ un conjunto con dos tipos de axiomas, lógicos y específicos de una $L$-teoría.
    \item $RDec$ un conjunto con dos reglas de deducción. \textit{Modus Ponens} (deducir $G$ de $F \then G$ y $F$), y \textit{Generalización} (deducir $\forall x\ F$ de $F$).
\end{itemize}

\begin{dfn}[Axiomas lógicos]$ $
    \begin{enumerate}
        \item[(I)] \textsc{Tautologías}
        \item[(II)] \textsc{Axiomas de cuantificadores}
        \begin{itemize}
            \item[(2.1)] Axioma de $\forall$:
            $$
                \forall v\ (F \then G) \then (F \then \forall v\ G)
            $$
            con $F, G \in \for(L)$ y $v \in V$ que no aparece libre.
            \item[(2.2)] Axioma de sustitución:
            $$
                \forall v\ F \then F(\quo{t}{v})
            $$
            con $F \in \for(L)$, $t \in \ter(L)$, $v \in V$.
        \end{itemize}
        \item[(III)] \textsc{Axiomas de igualdad}
        \begin{itemize}
            \item[(3.1)] Variables:
            $$
                x = x,\ x\in V
            $$
            \item[(3.2)] Funciones:
            $$
                x_1 = y_1 \then \cdots \then x_n = y_n \then f x_1 \ldots x_n = f y_1 \ldots y_n
            $$
            con $x_i, y_i \in V$, $f$ un símbolo de función $n$-aria.
            \item[(3.3)] Relaciones:
            $$
                x_1 = y_1 \then \cdots \then x_m = y_m \then R x_1 \ldots x_m = R y_1 \ldots y_m
            $$
            con $x_i, y_i \in V$, $R$ un símbolo de relación $m$-aria.
        \end{itemize}
    \end{enumerate}
\end{dfn}

\begin{pro}[Consecuencia semántica y consecuencia sintáctica]
    Sea $L$ un lenguaje, $F \in \for(L)$ un enunciado y $T$ una $L$-teoría: $T \vdash F \iff T \vDash F$.
\end{pro}

\section{Demostración formal}
\begin{dfn}[Demostración]
    Una \textbf{demostración} de $F \in \for(L)$ en una $L$-teoría $T$ es una secuencia finita de fórmulas de $L$ ($F_1, \ldots, F_n$) tal que $F_n$ es $F$ y donde para cada $i$ con $1 \leq i \leq n$:
    \begin{itemize}
        \item $F_i$ es un axioma lógico, ó
        \item $F_i$ se obtiene de $F_j, F_k$ donde $F_k$ es $F_j \then F_i$ mediante \textit{Modus Ponens}, ó
        \item $F_i$ se obtiene de $F_j$ mediante \textit{Generalización}, es decir, $F_i$ es $\forall v\ F_j$ para $v \in V$.
    \end{itemize}
\end{dfn}

\begin{dfn}[Teorema. Consecuencia sintáctica]
    Un \textbf{teorema} de una $L$-teoría $T$ es una fórmula $F \in \for(L)$, para la que existe una demostración de $F$ en $T$. Decimos que $F$ es un \textbf{teorema} de $T$, o $F$ es una \textbf{consecuencia sintática} de $T$. ($T \vdash F$).
\end{dfn}

\begin{obs}
    Si $(F_1, \ldots, F_n)$ es una demostración de $F$ en $T$, entonces $(F_1, \ldots, F_i)$ es una demostración de $F_i$ en $T$.
\end{obs}

\begin{dfn}[Longitud de una demostración]
    Si $(F_1, \ldots, F_n)$ es una demostración de $F$ en $T$, $n$ es la \textbf{longitud de la demostración}.
\end{dfn}

\begin{obs}$ $
    \begin{itemize}
        \item Los axiomas lógicos y los elementos de $T$ tienen demostraciones de longitud $1$.
        \item Si $T = \varnothing$ y $T \vdash F$ se dice que $F$ es un teorema del lenguaje $L$.
    \end{itemize}
\end{obs}

%% TODO: Añadir ejemplos de demostraciones (MUCHOS), empiezan en la página 23, día 10/10/2019

\begin{obs}[Ayudas para demostraciones]$ $
    \begin{itemize}
        \item Para demostrar $H_1 \bicond H_2$ basta demostrar $H_1 \then H_2$ y $H_2 \then H_1$. Luego se obtiene la fórmula original por tautologías.
        \item Para demostrar $H_1 \then \forall x\ H_2$ (con $x$ no libre en $H_1$) basta demostrar $H_1 \then H_2$. Se obtiene con generalización y el axioma de $\forall$.
        \item Para demostrar $\forall x \forall y\ F \bicond \forall y \forall x$ basta demostrar: $\forall x \forall y\ F \then \forall x\ F$ y $\forall y \forall x\ F \then \forall y\ F$.
    \end{itemize}
\end{obs}

\begin{dfn}[Demostración abreviada]
    Una \textbf{demostración abreviada} de $F \in \for(L)$ en $T$ tiene la misma definición que \textit{demostración} excepto que ahora consideramos las consecuencias semánticas de $T$ como parte de $T$.
\end{dfn}

\begin{obs}
    Una demostración abreviada de $F$ en $T$ es una demostración de $F$ en $T' = T \bigcup T_0$ donde $T_0 \subseteq \sdf{G \in \for(L) \mid T \vdash G}$.\\

    Dada una demostración abreviada de $F$ en $T$, se puede dar una demostración de $F$ en $T$. Sustituyendo en las líneas de la demostración de $F$ en $T \bigcup T_0$ cada $F_i$ teorema de $T$ por una demostración de $F_i$.
\end{obs}

\begin{pro}[$=$ como relación de equivalencia]\label{pro:2.1}
    Las fórmulas que expresan que $=$ es una relación de equivalencia son teoremas del lenguaje.
    \begin{enumerate}
        \item $\vdash x = x$.
        \item $\vdash x = y \then y = x$.
        \item $\vdash x = y \land y = z \then x = z$.
    \end{enumerate}
\end{pro}
\begin{proof}
    A completar. %% TODO: Página 27, día 15/10
\end{proof}

\section{Coherencia de teorías}
\begin{dfn}[Lenguaje coherente. Lenguaje incoherente.]
    Sea $L$ un lenguaje, $T$ una $L$-teoría. Decimos que $T$ es \textbf{incoherente} si existe $F \in \for(L)$ tal que $T \vdash F \land \neg F$. $T$ es \textbf{coherente} si no es \textit{incoherente}.
\end{dfn}

\begin{pro}[Caracterización de teorías incoherentes]\label{pro:2.3}
    Sea $L$ un lenguaje, $T$ una $L-teoría$. Las siguientes propiedades son equivalentes:
    \begin{enumerate}
        \item $T$ es incoherente.
        \item Existe $F \in \for(L)$ tal que $T \vdash \neg F$ y $T \vdash F$.
        \item Para toda $F \in \for(L)$ $T \vdash F$.
    \end{enumerate}
\end{pro}
\begin{proof}$ $
    \begin{itemize}
        \item[$(1 \implies 2)$] Sea $T$ incoherente, por definición, $\exists F \in \for(L)$ tal que $T \vdash F \land \neg F$.
        Entonces:
        \begin{align*}
            &F1:\ F \land \neg F & (T \vdash)\\
            &F2:\ F \land \neg F \then F & (taut.)\\
            &F3:\ F \land \neg F \then \neg F & (taut.)\\
            &F4:\ F & (MP(2, 1))\\
            &F5:\ \neg F & (MP(3, 1))
        \end{align*}
    \end{itemize}
    \item[$(2 \implies 3)$] Sea $G \in \for(L)$
        \begin{align*}
            &F1:\ F \then \neg F \then G & (taut.)\\
            &F2:\ F                      & (T \vdash)\\
            &F3:\ \neg F                 & (T \vdash)\\
            &F4:\ \neg F \then G         & (MP(1, 2))\\
            &F5:\ G                      & (MP(4, 3))
        \end{align*}
    \item[$(3 \implies 2)$] Partimos de $\forall F \in \for(L)$, $T \vdash F$. Como $F \land \neg F \in \for(L)$ entonces $T \vdash F \land \neg F$.
\end{proof}

% !TeX root = ../logica-matematica.tex

\chapter{Teoremas de Finitud y de la Deducci\'on.}

\section{Teorema de finitud}

\begin{thm}[Teorema de finitud]\label{thm:tf}
    Sea $L$ un lenguaje, $T$ una $L$-teoría y $F \in \for(L)$ una fórmula. Si $T \vdash F$ entonces existe $T_0 \subseteq T$ finito tal que $T_0 \vdash F$.
\end{thm}

\begin{cor}[Corolario primero al teorema \ref{thm:tf}]
    Si para $T_0 \subseteq T$, $T_0$ es coherente, entonces $T$ es coherente.
\end{cor}

\begin{cor}[Corolario segundo al teorema \ref{thm:tf}. Unión de teorías coherentes]
    Sea $L$ un lenguaje, sea $\sdf{T_i, i\in I}$ una familia de $L$-teorías, tal que $\forall i, j\in I$, $T_i \subseteq T_j$ o $T_j \subseteq T_i$. Si $T_i$ es coherente, $\forall i \in I$, entonces $\bigcup_{i\in I} T_i$ es una $L$-teoría coherente.
\end{cor}

\section{Teorema de la deducci\'on}

\begin{thm}[Teorema de la deducción]
    Sea $L$ un lenguaje, $T$ una $L$-teoría, $F$ un enunciado de $L$ y $G \in \for(L)$. Si $T \bigcup \sdf{F} \vdash G$ entonces $T \vdash F \then G$.
\end{thm}

\begin{proof}
    A completar. %%TODO: Página 29, día 17/10/2019
\end{proof}

% !TeX root = ../logica-matematica.tex

\chapter{Teorema de Igualdad.}


\section{Teorema de igualdad}

En esta sección vamos a ver algunos resultados acerca de igualdades. Vamos a apoyarnos en la proposición \ref{pro:2.1} por lo que convendría revisarla.

\begin{lm}[$=$ en las consecuencias sintácticas]\label{lm:igualsintactico}
    Sea $L$ un lenguaje, $t_i \in \ter(L)$ con $i = 1, 2, 3$, entonces:
    \begin{gather*}
        \vdash t_1 = t_1\\
        \vdash t_1 = t_2 \then t_2 = t_1\\
        \vdash t_1 = t_2 \land t_2 = t_3 \then t_1 = t_3
    \end{gather*}
\end{lm}

\begin{proof}
    Vamos a proceder con una demostración abreviada haciendo uso de la proposición \ref{pro:2.1}.
    \begin{align}
        &x=x &\text{Axioma de igualdad 3.1}\\
        &\forall x\ x=x&\text{Gen(1)}\\
        &\forall x\ x=x \then (x=x)(\quo{t_1}{x})&\text{Axioma de sustitución}\\
        &(x=x)(\quo{t_1}{x}): t_1=t_1&\text{MP(3, 1)}
    \end{align}
    \begin{align}
        &x=y \then y=x   &\vdash x=y \then y=x\\
        &\forall x\ F_1   &   \text{Gen(1)}\\
        &\forall x\ F_1 \then (x=y \then y=x)(\quo{t_1}{x})   &   \text{Axioma de sustitución}\\
        &t_1 = y \then y = t_1   &   \text{MP(3,2)}\\
        &\forall y\ F_4   &   \text{Gen(4)}\\
        &\forall y\ F_4 \then F_4(\quo{t_2}{y})   &   \text{Axioma de sustitución}\\
        &t_1 = t_2 \then t_2 = t_1   &   \text{MP(6, 5)}
    \end{align}
    \begin{align} %TODO: Demostrar
        &x=y \then y=x   &\text{se demuestra de forma análoga}
    \end{align}
\end{proof}

\begin{thm}[Teorema de igualdad]\label{thm:igualdad}
    Sea $L$ un lenguaje, $t, s_1, s_2 \in \ter(L)$, $x \in V$, $F \in \for(L)$ tal que $x$ es sustituible por $s_i$ en $F$. Sea $T$ una $L$-teoría. Entonces:
    \begin{enumerate}[(a)]
        \item $T \vdash s_1 = s_2 \implies T \vdash t(\quo{s_1}{x}) = t(\quo{s_2}{x})$
        \item $T \vdash s_1 = s_2 \implies T \vdash F(\quo{s_1}{x}) = F(\quo{s_2}{x})$
    \end{enumerate}
\end{thm}

\begin{proof}$ $
    \begin{enumerate}[(a)]
        \item Se da una demostración abreviada de $t(\quo{s_1}{x}) = t(\quo{s_2}{x})$ en $T$. Por inducción sobre la complejidad de $t$.\\\\

            Si $t\in V$:\\
                \begin{itemize}
                    \item $t$ es $x \implies t(\quo{s_i}{x})$ es $s_i$, y por tanto $T \vdash s_1 = s_2$ por hipótesis.
                    \item $t$ es $v \neq x \implies t(\quo{s_1}{x})$ es $t$, y por tanto $T \vdash t = t$ por el lema \ref{lm:igualsintactico}.
                    \item $t$ es $c$ se demuestra igual que el caso $v \neq x$
                \end{itemize}
            Si $t$ es $f t_1 \ldots t_m$. Por hipótesis de inducción se cumple que $T \vdash t_i(\quo{s_1}{x}) = t_i(\quo{s_2}{x})$, vamos a dar una demostración abreviada.\\
            \begin{align*}
                & t_1(\quo{s_1}{x}) = t_1(\quo{s_2}{x})                                                                                  &\ T\vdash \tag{$F_1$}\\
                & \vdots                                                                                                                 & \tag{$\cdots$}\\
                & t_m(\quo{s_1}{x}) = t_m(\quo{s_2}{x})                                                                                  &\ T\vdash  \tag{$F_m$}\\
                & x_1 = y_1 \then \cdots \then x_m = y_m \then f x_1 \ldots x_m = f y_1 \ldots y_m                                       & \text{Axioma de igualdad 3.2} \tag{$F_{m+1}$}\\
                & \forall x_1\ F_{m+1}                                                                                                   & \text{Gen($F_{m+1}$)} \tag{$F_{m+2}$}\\
                & \forall x_1\ F_{m+1} \then F_{m+1}(\quo{t_1(\quo{s_1}{x})}{x_1})                                                       & \text{Axioma de sustitución} \tag{$F_{m+3}$}\\
                & F_{m+1}(\quo{t_1(\quo{s_1}{x})}{x_1})                                                                                  & \text{MP($m+3$, $m+2$)}\tag{$F_{m+4}$}\\
                & \forall y_1\ F_{m+4}                                                                                                   & \text{Gen($F_{m+4}$)} \tag{$F_{m+5}$}\\
                & \forall y_1\ F_{m+4} \then F_{m+4}(\quo{t_1(\quo{s_2}{y})}{y_1})                                                       & \text{Axioma de sustitución} \tag{$F_{m+6}$}\\
                & F_{m+4}(\quo{t_1(\quo{s_2}{y})}{y_1})                                                                                  & \text{MP($m+6$, $m+5$)}\tag{$F_{m+7}$}\\
                & x_2 = y_2 \then \cdots \then x_m = y_m \then & \text{MP($m+7$, $1$)}\tag{$F_{m+8}$}\\
                &\then f t_1(\quo{s_1}{x}) x_2 \ldots x_m = t_1(\quo{s_2}{y}) f y_2 \ldots y_m &\\
                &\text{repetir el proceso m veces}&
            \end{align*}
        \item Demostramos por inducción sobre la complejidad de $F$. Como $F$ es $R t_1 \ldots t_m$ siendo $R$ un símbolo de relación $m$-aria (incluyendo $=$), basta demostrar:
            \begin{equation*}
                \tag{$+$} T \vdash F(\quo{s_1}{x}) \then F(\quo{s_2}{x})
            \end{equation*}
            ya que $T \vdash s_1 = s_2 \implies T \vdash s_2 = s_1$. De $(+)$ se obtiene:
            \begin{equation*}
                \tag{$++$} T \vdash F (\quo{s_2}{x}) \then F(\quo{s_1}{x})
            \end{equation*}
            y finalmente de $(+)$ y $(++)$ acabamos por obtener:
            $$
                T \vdash F(\quo{s_1}{x}) \bicond F(\quo{s_2}{x})
            $$
            por tautologías.
            \begin{itemize}
                \item Si $F$ es atómica, la demostración es análoga al apartado (a) usando el axioma de igualdad 3.3 en vez de 3.2.
                \item Si $F$ es $\neg G$.\\
                      Por hipótesis de inducción $T \vdash s_1 = s_2 \implies T \vdash G(\quo{s_1}{x}) \bicond G(\quox{s_2})$. De aquí podemos hallar $T \vdash \neg G(\quox{s_1}) \bicond \neg G(\quox{s_2})$ por tautologías.
                      Como $\neg(G(\quox{s_1}))$ es $(\neg G)(\quox{s_1})$ hemos terminado.
                \item Si $F$ es $G \then H$, por hipótesis de inducción y observando que $(G \then H)(\quox{s_i})$ es $(G(\quox{s_i}) \then H(\quox{s_i}))$ habríamos terminado.
                \item Si $F$ es $\forall y G$ distinguimos dos casos. Si $x$ no aparece libre, entonces $F(\quox{s_i})$ es $F$ y $T \vdash F \bicond F$ es una tautología.\\
                      Si $x$ aparece libre, entonces $x \neg y$ y como $x$ es sustituible por $s_i$ con $i = 1,2$ también tenemos que $y$ no aparece en $s_i$.\\
                      Además $(\forall y G)(\quox{s_i})$ es $\forall y (G(\quox{s_i}))$. Por hipótesis de inducción, $T \vdash s_1 = s_2 \implies T \vdash G(\quox{s_1}) \bicond G(\quox{s_2})$. Y como ya vimos en clase, podemos generalizar a ambos lados del bicondicional, con lo que obtenemos:
                      $$
                            T \vdash \forall y (G\quox{s_1}) \bicond \forall y (G\quox{s_2})
                      $$
                      es decir,
                      $$
                            T \vdash F \bicond F
                      $$
                      que es lo que queríamos demostrar
            \end{itemize}
    \end{enumerate}
\end{proof}

\begin{cor}[al teorema de igualdad (teorema \ref{thm:igualdad})]
    Con las mismas hipótesis que el teorema:
        \begin{enumerate}[(a')]
            \item $T \vdash s_1 = s_2 \then t(\quox{s_1}) \then t(\quox{s_2})$
            \item $T \vdash s_1 = s_2 \then F(\quox{s_1}) \then F(\quox{s_2})$
        \end{enumerate}
\end{cor}

\begin{proof}
    Si $s_1 = s_2$ fuera un enunciado, (a') y (b') se obtienen de (a) y (b) usando el teorema de la deducción (teorema \ref{thm:td}). En general, si $s_i$ es $s_i(\sucn{x})$, construimos $s_i'$ como $s_i'(\sucn{c})$ sustituyendo cada $x_j$ por el $c_j$ correspondiente.\\
    Por la proposición \ref{pro:2.3}, sea $L' = L \cup \sdf{\sucn{c}}$ y $s_i' \in \ter(L')$, tenemos que:
    \begin{itemize}
        \item[] (a')$\iff T \vdash s'_1 = s'_2 \then t(\quox{s'_1}) \then t(\quox{s'_2})$
        \item[] (b')$\iff T \vdash s'_1 = s'_2 \then F(\quox{s'_1}) \then F(\quox{s'_2})$
    \end{itemize}
    También sabemos que $T \vdash s_1 = s_2 \iff T \vdash s_1' = s_2'$. Como ahora son enunciados basta aplicar el teorema \ref{thm:td} y obtenemos el resultado que queremos demostrar.\\
    Además, también hemos demostrado que:
        \begin{gather*}
            T \vdash t(\quox{s_1}) = t(\quox{s_2}) \implies T \vdash t(\quox{s_1'}) = t(\quox{s_2'})\\
            T \vdash F(\quox{s_1}) = F(\quox{s_2}) \implies T \vdash F(\quox{s_1'}) = F(\quox{s_2'})
        \end{gather*}
\end{proof}

% !TeX root = ../logica-matematica.tex

\chapter{Teorema de Validez.}

\section{Introducción}

A lo largo del curso hemos estado interesados en las consecuencias semánticas, donde veíamos que $T \vDash F$ representaba que un enunciado $F$ se satisfacía en todos los modelos de $T$. A continuación, introdujimos la relación sintáctica ($T \vdash F$) como una versión más potente de la consecuencia semántica ($\vdash F \implies \vDash F$). Sin embargo, tenemos que comprobar que lo que hemos introducido es válido, es decir, que los teoremas son fórmulas válidas.

En esta sección vamos a ver las herramientas necesarias para demostrar que si $F$ es un teorema de $T$ ($T \vdash F$), es decir, existe una demostración de $F$ en $T$, entonces $F$ se satisface en $T$ ($T \vDash F$), es decir, $F$ se satisface en todos los modelos $\mtc A$ de $T$. Esto se enuncia formalmente con el teorema de validez.

\begin{thm}[Teorema de validez]\label{thm:tv}
    Sea $L$ un lenguaje, $T$ una $L$-teoría y $F \in \for(L)$ un enunciado.
    \begin{center}
        Si $T \vdash F \implies T \vDash F$. Informalmente:\\
        \textit{Los teoremas de $T$, que sean enunciados, se satisfacen en todos los modelos de $T$}.
    \end{center}
\end{thm}
\begin{proof}
    Más adelante. \hyperref[proof:tv]{Ir a la demostración}.
\end{proof}

Recordemos que una lógica consta de dos partes:
\begin{itemize}
    \item Una parte \textbf{sintáctica}: axiomas y teoremas ($T \vdash F$).
    \item Una parte \textbf{semántica}: estructuras y relaciones de satisfacción ($\mtc A \vDash F$, $T \vDash F$)
\end{itemize}

\begin{dfn}[Validez de una lógica]
    Decimos que una lógica es \textbf{válida} si satisface el teorema de validez (\ref{thm:tv}).
\end{dfn}

\begin{dfn}[Completitud de una lógica]
    Decimos que una lógica es \textbf{completa} si satisface el teorema de validez (\ref{thm:tv}) y su recíproco, es decir:
    $$
        T \vdash F \iff T \vDash F
    $$
\end{dfn}
\begin{obs}[Conjuntos de enunciados en una lógica completa]
    Si una lógica es completa, entonces sea $F$ un enunciado:
    $$
        \sdf{F \mid T \vdash F} = \sdf{F \mid T \vDash F}
    $$
\end{obs}

\begin{lm}[Validez de los axiomas lógicos]\label{lm:axval}
    Sea $L$ un lenguaje, los axiomas lógicos de $L$ son fórmulas válidas.
\end{lm}
\begin{proof}$ $
    \begin{enumerate}[(I)]
        \item \textsc{Tautologías}\\
        Ya lo vimos en la proposición \ref{pro:1.3}.
        \item \textsc{Axiomas de cuantificadores}
        \begin{itemize}
            \item[(2.1)] Axioma de $\forall$\\
                Sean $F, G \in \for(L)$ fórmulas, $F$ es $F(\sucn{x})$ y $G$ es $G(\sucn{x})$. Si $x$ no está libre en $G$:
                $$
                    H:\ \forall x (G \then F) \then G \then \forall x F
                $$
                Tenemos que ver que $H$ se satisface en cualquier estructura. Queremos llegar a:
                \begin{equation*}
                    \mtc A \vDash \forall x_1 \ldots \forall x_n H \text{ para toda $L$-estructura } \mtc A \tag{$\star$}
                \end{equation*}
                \begin{align*}
                    (\star) &\iff \text{ para todo } \sucn{a} \in A \implies \mtc A \vDash H(\sucn{a})\\
                            &\iff \mtc A \vDash (\forall x G \then F) (\sucn{a}) \tag{$1$} \\
                            &\iff \mtc A \vDash (G \then \forall x F) (\sucn{a}) \tag{$+$} \\
                    (+)     &\iff \text{Si } \mtc A \vDash G(\sucn{a}) \text{ entonces } \mtc A \vDash (\forall x F)(\sucn{a})
                \end{align*}
                Supongamos entonces:
                \begin{equation*}
                    \mtc A \vDash G(\sucn{a}) \tag{$2$}
                \end{equation*}
                y veamos que llegamos a
                \begin{equation*}
                    \mtc A \vDash (\forall x F)(\sucn{a}) \tag{$+++$}
                \end{equation*}
                Sea $x \neq x_i$, $(+++) \implies \forall b \in A$, $\mtc A \vDash F(b, \sucn{a})$. Llamemos $(3)$ a $\forall b \in A$, $\mtc A \vDash F(b, \sucn{a})$.\\
                Fijamos $b \in A$, por $(1)$, para este $b$ sabemos $\mtc A \vDash (G \then F)(b, \sucn{a})$.\\
                Por $(2)$, $\mtc A \vDash G(\sucn{a}) \implies \mtc A \vDash G(b, \sucn{a})$. Finalmente, como $x$ no aparece libre en $G$ podemos concluir que $(3)$ es cierto y por tanto, $(2) \implies (+++)$ con lo que se satisfacen $(+)$ y $(\star)$.

            \item[(2.2)] Axioma de sustitución\\
                Sea $F \in \for(L)$, $x \in V$m $t \in \ter(L)$ y $x$ sustituible por $t$ en $F$. Sea $\mtc A$ una $L$-estructura y $\sucn{a} \in A$.\\
                Veamos que si $\mtc A \vDash (\forall x F)(\sucn{a})$ entonces $\mtc A \vDash F(\quox{t})(\sucn{a})$.\\
                Por la proposición \ref{pro:1.2}, $\mtc A \vDash F(\quox{t})(\sucn{a}) \iff \mtc A \vDash F(t^{\mtc A}(\sucn{a}), \sucn{a})$.\\
                Como $t^{\mtc A}(\sucn{a}) \in A$, por $(1)$ tenemos que $\mtc A \vDash F(b, \sucn{a})$. En particular, también para $b = t^{\mtc A}(\sucn{a})$.
        \end{itemize}
        \item \textsc{Axiomas de igualdad}
        \begin{itemize}
            \item[(3.1)] Variables\\
            $$
                \mtc A \vDash (x=x)(a)\ \forall a \in A,\ \text{ entonces para todo } a \in A,\ a = a
            $$
            \item[(3.2)] Funciones\\
            $$
                \mtc A \vDash \forall x_1 \ldots \forall x_n \forall y_1 \ldots \forall y_n\ \ x_1 = y_1 \then \cdots \then x_n = y_n \then f x_1 \ldots x_n = f y_1 \ldots y_n
            $$
            Entonces es claro que: $\forall \sucn{a} \in A,\ \forall \sucn{b} \in A$ si $a_i = b_i$ entonces $f a_1 \ldots a_n = f b_1 \ldots b_n$.
            \item[(3.3)] Relaciones\\
            $$
                \mtc A \vDash \forall x_1 \ldots \forall x_m \forall y_1 \ldots \forall y_m\ \ x_1 = y_1 \then \cdots \then x_m = y_m \then R x_1 \ldots x_m = R y_1 \ldots y_m
            $$
            Entonces es claro que: $\forall \suc{b}{m} \in A,\ \forall \suc{b}{m} \in A$ si $a_i = b_i$ entonces $(\suc{a}{m}) \in R^{\mtc A}$ y $(\suc{b}{m} \in R^{\mtc A}) \implies (\sucn{a}) = (\suc{b}{m})$.
        \end{itemize}
    \end{enumerate}
\end{proof}

\section{Teorema de validez}

Una vez hemos demostrado que los axiomas son fórmulas válidas ya disponemos de las herramientas necesarias para demostrar el teorema de validez (teorema \ref{thm:tv}).

\begin{proof}[Demostración al teorema de validez]\label{proof:tv}
    La demostración sigue por inducción en la longitud de una demostración de $F$ en $T$.
    \begin{itemize}
        \item[$(n=1)$] Si la longitud de la demostración es $1$ o $F \in T$ o $F$ es un axioma lógico.\\
        Si $F \in T \implies \mtc A \vDash F,\ \forall \mtc A \vDash T$ y por tanto $T \vDash F$.
        Si $F$ es un axioma lógico, por el lema \ref{lm:axval} sabemos que $ \vDash F$ pues $F$ es una fórmula válida. Entonces $\mtc A \vDash F \forall \mtc A \implies T \vDash F$.

        \item[$(n>1)$]
        Supongamos que $F_n$ es $F$ y $(\sucn{F})$ es una demostración de $F$ en $T$.
        \begin{itemize}
            \item Si $F_n \in T$ o $F_n$ es un axioma lógico (como en $n=1$) entonces $T \vDash F$.
            \item Si $F_n$ es $\forall x F_i$ (aplicando \textit{Generalización} a algún $F_i$), por hipótesis de inducción, $T \vDash F_i (\iff T \vDash \overline{F_i} \text{ cualquier cierre universal de }F_i)$. Como todo cierre universal de $F_i$ también es cierre universal de $\forall x F_i$ entonces $T \vDash \forall x F_i$.
            \item Si $F_n$ se obtiene de $F_i$ y $F_j: F_i \then F_n$ aplicando \textit{Modus Ponens}, con $i, j < n$, por hipótesis de inducción:
            \begin{equation}
                T \vDash F_i\\
                T \vDash F_i \then F_n
            \end{equation}
            Veamos que $T \vDash F_n$. Sean $F_n:\ F_n(\sucn{x})$ y $F_i:\ F_i(\sucn{x})$, y $\mtc A \vDash T$, vamos a ver que $\mtc A \vDash \forall x_1 \ldots \forall x_n\ F_n$.\\
            Sean $(\sucn{a})\in A$, por (1) tenemos $T \vDash \forall x_1 \ldots \forall x_n F_i$. Para este $(\sucn{a})$ se tiene $\mtc A \vDash F_i(\sucn{a})$.
            Análogamente $\mtc A \vDash (F_i \then F_n)(\sucn{a}) \iff \mtc A (\sucn{a})$ entonces $\mtc{A} \vDash F_n(\sucn{a})$.\\
            Como tenemos el antecedente $\mtc A (\sucn{a})$, y también $\mtc A \vDash (F_i \then F_n)(\sucn{a})$ por H.I, obtenemos $\mtc{A} \vDash F_n(\sucn{a})$.
        \end{itemize}
    \end{itemize}
\end{proof}

\begin{cor}[Corolario primero al teorema \ref{thm:tv}]\label{cor:tv1}
    Sea $L$ un lenguaje, $T$ una $L$-teoría. Si existe $\mtc A \vDash T$ entonces $T$ es coherente.
\end{cor}

\begin{proof}
    Si $T$ es incoherente entonces existe un enunciado $F \in \for(L)$ tal que $T \vdash F \land \neg F$.\\
    Por el \hyperref[thm:tv]{teorema de validez}, $T \vDash F \land \neg F \implies T$ no puede tener modelos. Esto demuestra el contrarrecíproco.
\end{proof}

\begin{cor}[Corolario segundo al teorema \ref{thm:tv}]\label{cor:tv2}
    El corolario \ref{cor:tv1} es equivalente al \hyperref[thm:tv]{teorema de validez}.
\end{cor}

\begin{proof}
    No se vió.
\end{proof}

\part{Completitud de la l\'ogica de primer orden}
% !TeX root = ../logica-matematica.tex

\chapter{Teorías completas y teorema de Lindembaun.}

\section{Introducción}

Hasta ahora habíamos visto que una teoría completa era aquella que cumplía el teorema de validez y su recíproco. Podemos enunciar así una primera forma del teorema de Completitud.

\begin{thm}[Teorema de completitud. Primera forma.]\label{thm:tc1}
    Sea $L$ un lenguaje, $T$ una $L$-teoría y $F \in \for(L)$ entonces $T \vDash F \implies T \vdash F$.
\end{thm}

Sin embargo, podemos enunciarlo de otra forma, que es el recíproco al corolario \ref{cor:tv1}, ya que según el corolario 2 al teorema de validez, son equivalentes.
\begin{thm}[Teorema de completitud. Segunda forma.]\label{thm:tc2}
    Sea $L$ un lenguaje y $T$ una $L$-teoría, si $T$ es coherente entonces $T$ tiene un modelo.
\end{thm}


\begin{proof}[Demostración de que las dos formas del teorema de completitud son equivalentes]
    Vamos a ver que la primera forma del TC $\iff$ la segunda forma del TC.
    \begin{itemize}
        \item[$\implies$] Vamos a ver que si $T$ no tiene modelos, entonces $T$ es incoherente.\\
        Sea $F \in \for(L)$, $F$ un enunciado entonces, como $T$ no tiene modelos, trivialmente $T \vDash F \land \neq F$. Por la primera forma (teorema \ref{thm:tc1}) $T \vDash F \land \neq F \implies T \vdash F \land \neq F$ y entonces $T$ es incoherente.

        \item[$\implied$] Vamos a ver que sea $F \in \for(L)$, si $T \cancel\vdash F$ entonces $T \cancel\vDash F$, o expresado de otra forma:
        $$
            T \cancel\vdash F \implies T \cancel\vDash G, \text{ con $G$ el cierre universal de F}
        $$
        Puesto que $G$ es un enunciado podemos hacer uso del \hyperref[cor:td1]{corolario 1 al TD}. Si $T \cancel\vdash G$ entonces $T \cup \sdf{\neq G}$ es coherente.\\
        Si es coherente, por la segunda forma (teorema \ref{thm:tc2}) $T \cup \sdf{\neq G}$ tiene un modelo. Sea $\mtc A \vDash T \cup \sdf{\neq G}$, entonces $\mtc A \vDash T$ y $\mtc A \cancel\vDash G \implies T \cancel\vDash G \implies T \cancel\vDash F$.
    \end{itemize}
\end{proof}
Sin embargo, la demostración del teorema se verá más adelante en el curso.


\section{Teorías completas}


Vamos a ver que para cualquier teoría podemos añadir un conjunto de axiomas que la haga completa. Por ejemplo podemos pasar de la teoría de grupos (que ya vimos que es incompleta)
a grupos abelianos divisibles y sin torsión. GADST es una teoría completa.

\begin{eg}[Teorías completas e incompletas]
    \begin{enumerate}[(1)]
        \item Teoría de grupos:\\
        La teorías de grupos no es completa. Si $T \vdash \forall x \forall y\ xy = yx \implies T \vDash \forall x \forall y\ xy=yx$ por el teorema de validez, sin embargo esto no es cierto. Además, su negación tampoco lo es. Por tanto la teoría de grupos no es completa (no hay modelos).
        \item $T = \sdf{F \land \neq F}$ con $F$ un enunciado:\\
        Es una teoría completa. Toda teoría incoherente es trivialmente completa.
        \item Teoría de una estructura:\\
        La teoría de una estructura es completa. Sea $\mtc A$ una $L$-estructura, $T = \te(\mtc A) = \sdf{ F \in \for(L) \mid F \text{ enunciado y } \mtc A \vDash F}$.\\
        Sea $F\in \for(L)$ un enunciado, entonces $A \vDash F$ o $A \vDash \neq F \implies F \in \te(A)$ o $\neq F \in \te(A)$ y por tanto $\te(A) \vdash F$ o $\te(A) \vdash \neq F$ con una demostración de longitud $1$.
    \end{enumerate}
\end{eg}

\section{Teorema de Lindebaum}

\part{Ap\'{e}ndices}

%\renewcommand{\chaptername}{Apéndice}
%\renewcommand{\thechapter}{\Alph{chapter}}
%\setcounter{chapter}{0}

%% !TeX root = ../apuntes-ea.tex

\chapter{Ejercicios}

\section{Hoja 1}

\begin{ex}[H1.2]
	Sean $a,b,c \in G = (-1,1)$. Probamos las propiedades de los grupos.
	\begin{itemize}
		\item \textbf{Asociatividad:}
		\begin{align*}
			(a \ast b) \ast c = \left(\frac{a+b}{1+ab}\right) \ast c = \frac{\left(\frac{a+b}{1+ab}\right) + c}{1 + \left(\frac{a+b}{1+ab}\right)c} = \frac{\frac{a+b+c+abc}{1+ab}}{\frac{1+ab+ac+bc}{1+ab}} = \frac{a+b+c+abc}{1+ab+ac+bc} \\
			a \ast (b \ast c) = a \ast \left(\frac{b+c}{1+bc}\right) = \frac{a+\left(\frac{b+c}{1+bc}\right)}{1+a\left(\frac{b+c}{1+bc}\right)} = \frac{\frac{a + b + c + abc}{1+bc}}{\frac{1+ab+ac+bc}{1+bc}} = \frac{a+b+c+abc}{1+ab+ac+bc}
		\end{align*}
		\item \textbf{Elemento neutro:} es el $0$ ya que $x\ast 0 = \frac{x+0}{1+x\cdot 0} = \frac{x}{1} = x$ y además $0 \ast x = \frac{0+x}{1+0\cdot x} = \frac{x}{1} = x$
		\item \textbf{Elemento inverso:} la ecuación
		\begin{align*}
			x \ast \inv{x} = 0 \iff \frac{x + \inv{x}}{1+x\inv{x}} = 0 \iff \inv{x} = -x
		\end{align*}
		siempre tiene solución y ocurre lo mismo para la ecuación $\inv{x} \ast x = 0 \iff \inv{x} = -x$
		\item \textbf{Clausura:} tenemos que probar que si $x,y \in (-1,1)$ entonces $x \ast y \in (-1,1)$. Consideramos $f(x,y) = x \ast y = \frac{x+y}{1+xy}$. Derivando tenemos que $\nabla f(x,y) = (\frac{1}{(1+xy)^2},\ \frac{1}{(1+xy)^2}) \neq 0, \forall x,y \in [-1,1]\times[-1,1]$. Si el máximo no se alcanza en ningún sitio de dentro del cuadrado $(-1,1)\times(-1,1)$ se tendrá que alcanzar en el borde.
		\begin{itemize}
			\item Fijado $x = 1$ tenemos que $f(1,y) = \frac{1+y}{1+y} = 1 \implies f(1, -1 < y < 1) < 1$ porque si $f(1, -1 < y < 1)$ tomara un valor mayor que $1$ habría un máximo en $(-1,1)\times(-1,1)$ y esto no puede ser pues $\nabla f$ no se anula en el cuadrado.
			\item Fijado $x = -1$ tenemos que $f(-1, y) = \frac{y-1}{1-y} = -1 \implies f(-1, -1 < y < 1) > -1$ por la misma razón que antes.
			\item Hacemos lo mismo fijando la $y$ y variando la $x$.
		\end{itemize}
		En el borde (que no está incluido) se alcanzan máximo y mínimo que acotan a $f$ en el cuadrado:
		\begin{align*}
			-1 < f(x,y) = x \ast y < 1,\quad \forall x, y \in G
		\end{align*}
	\end{itemize}
\end{ex}

\begin{ex}[H1.3]Hallar los inversos de los siguientes elementos, cada uno en su grupo correspondiente:
	\begin{enumerate}
		\item $o(\overline{11})$ en $\uds{\Z^\ast/23\Z}$ es $22$ porque $23 \cdot 22 \equiv 1 \mod 23$
		\item $o(\overline{5})$ en $\uds{\Z^\ast/31\Z}$ es $3$ porque $5 \cdot 3 \equiv 1 \mod 31$
	\end{enumerate}
\end{ex}

\begin{ex}[H1.33]
	\label{ex:h1.33}
	Sea $G$ un grupo. Suponed que existe un único $a \in G$ de orden 2. Demostrad que $a \in Z(G)$.
\end{ex}

\begin{proof}
	Recordamos que $a \in Z(G) \iff ga = ag,\ \forall g \in G$. Definimos el isomorfismo de conjugación $\phi_g (x) = gx\inv{g}$ para algún $g$. Como $\phi_g$ es isomorfismo lleva elementos de orden $n$ en elementos de orden $n$. Entonces $\phi_g(a) = a$ ya que $a$ es el único elemento de orden 2. Por tanto $ga\inv{g} = a \implies ga = ag \implies a \in Z(G)$.
\end{proof}

\section{Hoja 2}

\begin{ex}[H2.1]
	\label{ex:h2.1}
	Se considera el tercer grupo diédrico $D_3$. Se pide hallar lo siguiente:
	\begin{enumerate}
		\item Las clases de conjugación de cada uno de sus elementos.
		\begin{proof}
			Las clases dan una partición del grupo. Si un elemento pertenece a una clase, entonces la clase de ese elemento también es la clase a la que pertenece.
			\begin{itemize}
				\item $cl(e) = \{e\}$
				\item $cl(B)$? Sabemos que $|cl(B)| = [G:C(B)]$. Sabemos que $\gen{B} =\{1, B, B^2\} \subset C(B)$ luego $|C(B)| \geq 3$. Si hubiera más elementos en $C(B)$ tendríamos que $|C(B)| = 6$ pues $C(B) < D_3$. Esto no ocurre porque sabemos que $B$ no conmuta con todos los demás elementos. Por ejemplo $BA \neq AB$. Por tanto $|C(B)| = 3 \implies |cl(B)| = [D_3:C(B)]= 6/3 = 2$. Es claro que $B \in cl(B)$. Además, como $cl(B)$ contiene elementos transformados por el isomorfismo conjugación sabemos que el otro elemento que hay tiene orden 3. El único elemento que queda de orden 3 es $B^2 \implies cl(B) = \{B, B^2\}$.
				
				\item $cl(A)?$ Sabemos que $A$ no conmuta con todos ($A \not\in Z(D_3)$) luego $|C(A)| < 6$. Sabemos que $\gen{A} = \{1, A\} < C(A)$. Además, como $C(A)$ es un (sub)grupo sabemos que no puede haber más elementos porque si los hubiera, $|\gen{A}| \divides |C(A)| \implies C(A) \geq 6$ pero ya hemos visto que no puede ser. Es decir que $|cl(A)| = [D_3:D(A)] = 6 / 2 = 3$. Por tanto $cl(A)$ incluye los 3 elementos que nos quedan: $cl(A) = \{A, AB, AB^2\}$.
			\end{itemize}
		\end{proof}
	
		\item Los elementos de $\text{Int}(D_3)$.
		\item Los centralizadores $C_{D_3}(x)$ para cada $x \in D_3$
		\item Los normalizadores $N(H)$ para cada $H < D_3$.
	\end{enumerate}
\end{ex}

\begin{ex}[H2.2]
	
	\begin{proof}
		Obtenidas las clases en el ejercicio \nameref{ex:h2.1} se verifica que $|D_3| = |cl(e)| + |cl(B)| + |cl(A)| = 1 + 2 + 3 = 6$
	\end{proof}
	
\end{ex}

\begin{ex}[H2.6]
	Sea $G$ un grupo. ¿Verdadero o falso?
	\begin{enumerate}
		\item $H < G$ y $H$ conmutativo implica $H \normsub G$.
		\item $H < G$ y $|H| = 2$ implica $H \normsub G$.
		\item Si $\varphi: G \to G_1$ es un homomorfismo de grupos, entonces $\ima \varphi \normsub G$
		\item Si $H \normsub K$ y $K \normsub G$ entonces $H \normsub G$
		\item Si $H \normsub G$ y $|H| = m$ entonces $H$ es el único subgrupo de $G$ de orden $m$.
		\item Si $H \normsub G$ entonces $H < Z(G)$.
		\begin{proof}[FALSO]
			Contraejemplo: En $G = D_4$ tomamos $H = \gen{B^2} = \{1, B, B^2, B^3\} \not\subset Z(D_4) = \{1, B^2\}$.
		\end{proof}
	\end{enumerate}
\end{ex}

\begin{ex}[H2.10]
	\begin{proof}
		Fijado $n$ y definida $\alpha_n : G \to G,\ x \mapsto x^n$ tenemos que $\alpha_n$ es un homomorfismo de grupos. Además podemos expresar $H_2 = \ker \alpha_n \implies H_2 \normsub G$. Además también tenemos que $H_1 = \ima \alpha_n < G$. Veamos que $H_1 \normsub G$. Es decir, que $gH_1\inv{g} = H_1,\ \forall g \in G$. Para ello tomamos $x_1^n \in H_1$ y lo conjugamos $gx_1^n\inv{g} = (gx_1\inv{g})^n$ por ser $\alpha$ homomorfismo de grupos. En particular $(gx_1\inv{g})^n \in \ima \alpha \implies (gx_1\inv{g})^n \in H_1 \implies (gx_1\inv{g})^n = x_2^n$ para algún $x_2 \in H_1 \implies H_1 \normsub G$.
	\end{proof}
\end{ex}

\begin{ex}[H2.13] Si $A$ es un grupo abeliano con $n$ elementos y $k$ es un entero primo con $n$, demostrad que la aplicación $\varphi : A \to A$ definida por $\varphi(a) = a^k$ es un isomorfismo.
	\begin{itemize}
		\item $\varphi$ homomorfismo de grupos.
		\begin{proof}
			\begin{align*}
				\varphi(a)\varphi(b) = a^kb^k = (ab)^k = \varphi(ab)
			\end{align*}
		\end{proof}
		\item $\varphi$ biyectiva $\iff \varphi$ inyectiva ya que dominio y codominio coinciden
		\begin{proof}
			$\ker \varphi = \{a \in A \mid \varphi(a) = a^k = 1\}$. Probaremos que $a^k = 1 \iff a = 1$ y por tanto que $\ker \varphi = \{1\} \implies \varphi$ inyectiva. Sabemos que $a^k = 1 \iff o(a^k) = 1$. Sea $t = o(a) \divides n$. Distinguimos dos casos
			\begin{itemize}
				\item Si $t = 1$ entonces $a = 1$ y ya está
				\item Si $t > 1$ entonces $o(a^k) = \frac{t}{mcd(k, t)} = \frac{t}{1} > 1$ contradicción. Luego necesariamente $t = o(a) = 1$.
			\end{itemize}
		\end{proof}
	\end{itemize}
\end{ex}

\begin{ex}[H2.22]
	\label{ex:h2.22}
	Demostrad que si $G$ es un grupo no conmutativo y tiene orden $p^3$ ($p$ un número primo) entonces $Z(G)$ tiene orden $p$.
	
	\begin{proof}
		Sabemos que $Z(G) < G \implies |Z(G)| \divides |G| \implies |Z(G)| \in \{1, p, p^2, p^3\}$
		\begin{itemize}
			\item $|Z(G)| \neq p^3$ porque en tal caso $G$ sería conmutativo
			\item $|Z(G)| \neq 1$ porque $G$ es un p-grupo y por tanto su centro no es el trivial.
			\item Si $|Z(G)| = p^2$ entonces $|G/Z(G)| = p \implies G/Z(G)$ es cíclico lo que no es posible si $G$ no es abeliano.
		\end{itemize}
	
		Por descarte concluimos que $|Z(G)| = p$.
	\end{proof}
\end{ex}


\begin{ex}[H2.25]
	Sabemos que $\autom{\Z/12\Z} \isom \uds{\Z^\star/12\Z}$ donde $\Z^\star/12\Z$ es el grupo multiplicativo $(\{\overline{1}, \overline{2}, \overline{3}, \dots,  \overline{11}\}, \cdot)$. Queda $\uds{\Z^\star/12\Z} = (\{\overline{1}, \overline{5}, \overline{7}, \overline{11}\}, \cdot )$ y además da la casualidad que $\forall x \in \uds{\Z^\star/12\Z},\ o(x) = 2$ (todos los elementos son su propio inverso) por lo que no tenemos restricciones al definir $f: \Z/2\Z \to \autom{\Z/12\Z}$:
	\begin{align*}
		f:\Z/2\Z & \to \autom{\Z/12\Z} \isom \uds{\Z^\star/12\Z} \\
		e = \overline{0} &\mapsto 1 \\
		\overline{1} &\mapsto \{\overline{1}, \overline{5}, \overline{7}, \overline{11}\}
	\end{align*}
\end{ex}

\begin{ex}[H2.26]
	Sea $|G_1| = m, |G_2| = n,\ mcd(m, n) = 1$. Si $f:G_1 \to G_2$ es h. de g. sabemos que $o(f(a)) \divides o(a),\ \forall a \in G_1$. Además $o(a) \divides m \land o(f(a)) \divides n$ por el teorema de Lagrange (\ref{thm:lagrange}).
	\begin{align*}
		\begin{cases}
		o(a) \divides m \land o(f(a)) \divides n \\
		mcd(m, n) = 1 \\
		o(f(a)) \divides o(a)
		\end{cases} \implies o(a) = o(f(a)) = 1, \forall a \in G_1
	\end{align*}
	Por lo que solo puede haber un homomorfismo entre ellos y además es el trivial $f(a) = e_{G_2}$.
\end{ex}

\begin{ex}[H2.19]
	Definimos una función $f:[0, 2\pi] \subset \R \to \mathbb{S}^1, \alpha\mapsto \cos \alpha + i \sin \alpha$. Esta función tiene la propiedad de que $f(\alpha) \cdot f(\alpha') = \cos(\alpha + \alpha') + i\sin(\alpha + \alpha') = f(\alpha + \alpha')$ y por tanto es un h. de g.\footnote{A la izquierda (en $\R$) sumamos pero a la derecha (en $\mathbb{S}^1$) multiplicamos.} entre $R$ y $\mathbb{S}^1$.
	
	Un elemento de $\cos \alpha + i \sin \alpha \in \mathbb{S}^1$ es de torsión $\iff \exists n \mid (\cos \alpha + i \sin \alpha )^n = 1$. Ahora bien $(\cos \alpha + i \sin \alpha )^n = \cos n\alpha + i \sin n\alpha = 1 \iff n\alpha = k 2\pi$. 
\end{ex}

\section{Hoja 4}

\begin{ex}[H4.11] Hallar los subgrupos de Sylow de $S_5$.
	Sabemos que $|S_5| = 5! = 2^3\cdot3\cdot 5$ y por \nameref{thm:sylow1} tenemos lo siguiente:
	\begin{itemize}
		\item $\exists P_2,\ |P_2| = 2^3 = 8$.
		\item $\exists P_3,\ |P_3| = 3$. Además en $S_5$ hay $\binom{5}{3}2! =20 $ 3-ciclos y en cada $gP_3\inv{g} = \{1, a, a^2 \mid o(a) = 3\}$ hay 2 elementos de orden 3 distintos. Además, $g_1P_3\inv{g_1} \cap g_2P_3\inv{g_2} = \{e\}$ porque si su intersección fuera más grande entonces serían el mismo subgrupo (porque son cíclicos). Es por esto que tenemos que repartir 40 elementos dando 2 a cada 3-grupo con lo que obtenemos $n_3 = 20 / 2 = 10$ 3-subgrupos de Sylow en $S_5$.
		\item $\exists P_5,\ |P_5| = 5$. Además en $S_5$ hay $4! = 24$ 5-ciclos (elementos de orden 5) y en cada $gP_5\inv{g} = \{1, a, a^2, a^3, a^4 \mid o(a) = 5\}$ tenemos $4$ elementos de orden $5$ distintos. Además, $g_1P_5\inv{g_1} \cap g_2P_5\inv{g_2} = \{e\}$ porque si su intersección fuera más grande entonces serían el mismo. Así, tenemos $24$ 5-ciclos a repartir entre los diferentes $gP_5\inv{g}$ dando 4 5-ciclos a cada 1. Por tanto tenemos $n_5 = 24/4 = 6$ 5-subgrupos de Sylow en $S_5$.
	\end{itemize}
\end{ex}

\begin{ex}[H4.18]
	Demostrar que todo grupo de orden $|G| = 5^3 \cdot 7^3$ tiene un subgrupo normal de orden 125.
	
	\begin{proof}
		\nameref{thm:sylow1} $\implies \exists P_5 < G,\ |P_5| = 125$. \nameref{thm:sylow3} $\implies n_5 \divides 7^3 \land n_5 \equiv 1 \mod 5$ es decir $n_5 \in \{1, 7, 49, 343\} \land n_5 \in \{1, 6, 11, \dots\}$. Como ni $49$ ni $343$ son conguentes con $1$ módulo 5 tenemos que $n_5 = 1 \implies P_5 \normsub G$.
	\end{proof}
\end{ex}

\begin{ex}[H4.20] Hallad todos los grupos abelianos de órdenes 36, 64, 96 y 100.
	\begin{enumerate}
		\item $|G| = 36 = 2^2\cdot3^2$
		\begin{proof}
			\nameref{thm:sylow1} $\implies \exists P_2,\ |P_2| = 4 \land \exists P_3,\ |P_3| = 9$. Además $G$ abeliano $\implies P_2, P_3 \normsub G \implies G \isom P_2 \times P_3$. Estudiamos los grupos de orden 4 y de orden 9
			\begin{itemize}
				\item $|P_2| = 4$ entonces $P_2 \isom \Z/2\Z \times \Z/2\Z \lor P_2 \isom \Z/4\Z$
				\item $|P_3| = 9$ entonces $P_3 \isom \Z/3\Z \times \Z/3\Z \lor P_3 \isom \Z/9\Z$
			\end{itemize}
			Como $G \isom P_2 \times P_3$ tenemos 4 posibles grupos abelianos de orden 36.
		\end{proof}
	\end{enumerate}
\end{ex}

\begin{ex}[H4.22]
	Hallar todos los grupos abelianos de orden 175.
	
	\begin{proof}
		$|G| = 5^2\cdot 7$. Por el \nameref{thm:sylow1} tenemos que $\exists P_5, P_7 < G$ con $|P_5| = 25,\ |P_7| = 7$ y además por ser $G$ abeliano tenemos que $P_5, P_7 \normsub G \implies G \isom P_5 \times P_7$. Estudiamos los grupos de orden 25 y de orden 7:
		\begin{itemize}
			\item $|P_5| = 25 \land P_5$ abeliano $\implies P_5 \isom \Z/25\Z \lor P_5 \isom \Z/5\Z \times \Z/5\Z$. En ambos casos $P_5$ es producto directo de cíclicos pues $\Z/n\Z$ es cíclico.
			\item $|P_7| = 7 \land P_7$ abeliano $\implies P_7 \isom \Z/7\Z$. Ocurre lo mismo que con $P_5$.
		\end{itemize}
		Concluimos que $G \isom \Z/25\Z \times \Z/7\Z \lor G \isom \Z/5\Z \times \Z/5\Z \times \Z/7\Z$. Los dos casos son abelianos por ser producto directo de grupos cíclicos.
	\end{proof}
\end{ex}

\begin{ex}[H4.23]
	¿Cuántos elementos de orden 3 puede tener un grupo abeliano de orden 36?
	
	\begin{proof}
		$|G| = 36 = 2^2 3^2$. \nameref{thm:sylow1} $\implies \exists P_2, P_3 < G,\ |P_2| = 4,\ |P_3| = 9$. Estudiamos los grupos de órdenes 4 y 9:
		\begin{itemize}
			\item $|P_3| = 9 \implies P_3 \isom \Z/9\Z \lor P_3 \isom \Z/3\Z \times \Z/3\Z$
			% TODO: terminar
		\end{itemize}
	\end{proof}
\end{ex}

\chapter{\'{I}ndices}

%%TODO por qué no podemos utilizar only= para hacer whitelisting en lugar de blacklisting? que alguien lo arregle parfavar
\renewcommand{\listtheoremname}{Lista de definiciones}
\listoftheorems[ignore={thm,eg,pro,cor,obs,nota,lm,ex,th_ex}]

\renewcommand{\listtheoremname}{Lista de teoremas}
\listoftheorems[onlynamed,ignore={dfn,eg,cor,obs,nota,ex,th_ex}]

\renewcommand{\listtheoremname}{Lista de ejemplos}
\listoftheorems[onlynamed,ignore={dfn,thm,pro,cor,obs,nota,lm,ex,th_ex}]

\renewcommand{\listtheoremname}{Lista de ejercicios}
\listoftheorems[ignore={dfn,thm,pro,cor,obs,nota,lm,eg}]

\printglossaries

\bibliographystyle{alpha}

\end{document}
